% BEGIN INCLUDED SGA7I_English_current_buildable_source_20260823/components/150_expose_IX_lemma_7_2_1.tex
% Expose IX, 7.2 and Lemma 7.2.1 with proof 7.2.1.1--7.2.1.5.
% French transcription lines 1889--1938; authority scan physical pp. 414--418
% (scan indices 413--417; source folios 402--406).  All pages were checked
% directly in four overlapping 1100-dpi bands.  No unresolved feature needed
% higher-resolution escalation.
% Source notation corrected transparently: the adic base is Spf(V), not the
% printed Spf(A), since A is the varying group scheme; likewise Phi_o and
% Phi_n live over S_o and S_n, although 7.2.1.3 prints S.  The parenthetical
% source citation for representability in 7.2.1.3 is blank and is not guessed.

\medskip
\noindent
\SGARefTarget{sga7.sga7i.eix.n7-2}{7.2}. 现在记下如下引理.

\medskip
\noindent
\underline{引理 \SGARefTarget{sga7.sga7i.eix.n7-2-1}{7.2.1}.} 设 $S$ 为 诺特环 $V$ 的谱, 此环对于 $V$ 的理想 $\underline m$ 所定义的拓扑分离且完备. 令 $C$ 为 $S$ 上的群概形 $\Lambda$ 所成的范畴, 要求其在 $S$ 上光滑且有限型, $\Lambda/\Lambda^o$ 为 $S$ 上的有限平展群概形, 而 $\Lambda^o$ 为 阿贝尔概形 $B$ 被等平凡环面 $T$ 的扩张 (SGA~3 IX~1.1). 令 $C'$ 为 $\widehat S=\operatorname{Spf}(V)$ 上有限型形式概形范畴 (EGA~I~10.3.3) 中的群对象所成的范畴, 并考虑由 $C$ 到 $C'$ 的函子
\[
 A\longmapsto\widehat A
\]
它取沿 $A_o=A\times_S S_o$ 的形式完备化, 其中 $S_o=\operatorname{Spec}(V/\underline m)$.\footnote{在此陈述中的两处, 印本都写作 $\operatorname{Specf}(A)$, 在下文还写作 $\operatorname{Spec}(A/\underline m^{n+1})$. 由于 $V$ 是理想进底环, 而 $A$ 表示变动的群概形, 类型正确的底应为 $\operatorname{Spf}(V)$, 且 $S_n=\operatorname{Spec}(V/\underline m^{n+1})$.} 则:

\smallskip
\noindent
\underline{(a) 上述函子是全忠实的.}

\smallskip
\noindent
\underline{(b) 设 $\mathcal A$ 为 $C'$ 的一个对象, 在 $\widehat S$ 上光滑}; 即假设通常的概形
\[
 A_n=\mathcal A\times_{\widehat S}S_n,
 \qquad
 S_n=\operatorname{Spec}(V/\underline m^{n+1}),
\]
在 $S_n$ 上光滑. 假设 $A_o/A_o^o$ 为 $S_o$ 上的有限平展群概形, 而 $A_o^o$ 是 阿贝尔概形 $B_o$ 被等平凡环面 $T_o$ 的扩张. 置
\[
 \mathcal A^o=(A_n^o)_{n\geq0},
 \qquad
 \mathcal A/\mathcal A^o=(A_n/A_n^o)_{n\geq0}.
\]
\underline{则 $\mathcal A/\mathcal A^o$ 在 $\widehat S$ 上有限平展, 且 $\mathcal A^o$ 以本质唯一且关于 $\mathcal A$ 函子性的方式, 成为 $\widehat S$ 上的形式 阿贝尔概形 $\mathcal B$ 被 $\widehat S$ 上的形式环面 $\mathcal T$ 的扩张, 它们分别约化为 $B_o$ 和 $T_o$. 形式群概形 $\mathcal A$ 属于 $C\longrightarrow C'$ 的本质像, 或者我们也称它是「可代数化的」, 当且仅当 $\mathcal B$ 可代数化.}

\medskip
\noindent
\SGARefTarget{sga7.sga7i.eix.n7-2-1-1}{7.2.1.1}. 首先注意, 在任意概形 $S$ 上, 若 $A$ 是光滑群概形, 且 $A^o$ 是某个 阿贝尔概形被环面的扩张, 则这一扩张结构本质唯一, 且在显然的意义下关于 $A$ 为函子性的. 事实上, 该断言立即归结为从环面 $T$ 到 阿贝尔概形 $B$ 的每个态射均为零这一事实. 此事实在几何纤维上众所周知 [15], 继而由 SGA~3 IX~5.3 得到结论. 将该结果用于 \SGARefLink{sga7.sga7i.eix.n7-2-1}{7.2.1} 中 $S_n$ 上光滑群概形 $A_n$ 的相容系统, 就得到 $S$ 上形式群概形的类似结果. 此外, 在 \SGARefLink{sga7.sga7i.eix.n7-1}{7.1} 中已注意到, 若 $A_o^o$ 为 阿贝尔概形被环面的扩张, 则 $A_n^o$ 也如此. 这就证明了 (b) 的第一个断言.

\medskip
\noindent
\SGARefTarget{sga7.sga7i.eix.n7-2-1-2}{7.2.1.2}. 其次假设 $S$ 上的形式群概形 $\mathcal A$ 形如 $\widehat A$, 其中 $A\in\operatorname{ob}C$. 若 $A^o$ 为 阿贝尔概形 $B$ 被环面 $T$ 的扩张, 则显然 $\mathcal B\simeq\widehat B$. 因而 $\mathcal A$ 可代数化蕴涵 $\mathcal B$ 可代数化.

反之, 假设 $\mathcal B$ 可代数化, 即它来自某个 阿贝尔概形 $B$, 来证明 $\mathcal A$ 可代数化. 首先证明 $\mathcal A^o=(A_n^o)$ 可代数化. 由 VIII~\SGARefLink{sga7.sga7i.eviii.n3-7}{3.7}, $B_n$ 被 $T_n$ 的扩张 $A_n^o$ 可以解释为同态
\begin{equation}\tag{*}
 \underline M_n\longrightarrow B'_n,
\end{equation}
其中 $\underline M_n$ 是 $S_n$ 上与 $T_n$ 对偶的扭常值群 (SGA~3 X~5.7), 而 $B'_n$ 为 $B_n$ 的对偶 阿贝尔概形. 限制到 $S_o$ 或任意 $S_n$, 诱导从 $S$ 上的等平凡环面范畴到 $S_o$ 上相应范畴的等价 (SGA~3 X~3.2). 因而各 $T_n$ 来自 $S$ 上的一个环面 $T$, 它精确到唯一同构被确定. 将其对偶群记为 $\underline M$; 它在 $S_n$ 上的限制典范同构于 $\underline M_n$.

$T$ 的等平凡性意味着 $\underline M$ 是一族在 $S$ 上有限的 $S$-概形 $\underline M^i$ 的不交并 (SGA~3 X~5.11). 立即得到, 各同态 $(*)$ 来自唯一的同态
\begin{equation}\tag{**}
 \underline M\longrightarrow B',
\end{equation}
其中 $B'$ 为 $B$ 的对偶 阿贝尔概形. 而此同态又解释为 $B$ 被 $T$ 的一个扩张, 记作 $A^o$; 显然这给出 $\mathcal A^o$ 的一个代数化.

\medskip
\noindent
\SGARefTarget{sga7.sga7i.eix.n7-2-1-3}{7.2.1.3}. 另一方面, 根据假设, $A_o/A_o^o$ 是 $S_o$ 上的有限群概形 $\Phi_o$, 而各 $A_n/A_n^o$ 同样是 $S_n$ 上的有限群概形 $\Phi_n$; 它们必在 $S_n$ 上平展.\footnote{在这些句子中, 印本写作 $S$, 而不是纤维的底 $S_o$ 和 $S_n$, 并且用于推出 $\Phi_n$ 可表性的括号内文献引用留白. 此处根据引理的陈述恢复这些底; 不猜测缺失的文献引用.} 因而存在 $S$ 上的平展群概形 $\Phi$, 精确到唯一同构被确定, 其形式完备化为各 $\Phi_n$ 所定义的形式群概形 $\mathcal A/\mathcal A^o$.

由于 $G=A^o$ 可除, $S$ 上如 $\Phi$ 这样的有限局部常值群被群概形 $A^o$ 的 fppf 层扩张所成的范畴, 典范等价于以层
\[
 H=\underline{\operatorname{Hom}}_{\mathrm{gr}}(\Phi,G),
\]
为结构群的挠子范畴, 且此等价与任意基变换相容. 类似地, 将 $\Phi_n$ 被 $G_n$ 的扩张解释为以
\[
 H_n\simeq
 \underline{\operatorname{Hom}}_{\mathrm{gr}}(\Phi_n,G_n),
\]
为结构群的挠子, 则各 $A_n$ 定义了一个以各 $H_n$ 为结构群的挠子 $P_n$ 的相容系统. 由于 $G$ 为 阿贝尔概形被环面的扩张, 容易推出 $H$ 为\underline{在 $S$ 上有限且局部自由的}群概形. 因此, 函子
\[
 P\longmapsto\widehat P=(P_n)
\]
是 $S$ 上以 $H$ 为结构群的挠子范畴, 与以 $\widehat H$ 为结构群的「形式」挠子 $\widehat P$ 的范畴之间的等价; 后者即以各 $H_n$ 为结构群的挠子 $P_n$ 的相容系统.

因而, 各扩张 $A_n$ 所定义的 $P_n$ 来自 $S$ 上一个以 $H$ 为结构群的挠子 $P$. 得到 $\Phi$ 被 $G=A^o$ 的一个扩张 $A$, 它给出 $\mathcal A$ 的代数化, 只需再检验其可表性. 将 $A$ 看作 $\Phi$ 上以 $A^o_\Phi$ 为结构群的挠子. 对有限、满射、局部自由态射
\[
 P\times_S\Phi\longrightarrow\Phi,
\]
作基变换后, 此挠子平凡, 因为基变换 $P\longrightarrow S$ 使 $\Phi$ 被 $A^o$ 的扩张 $A$ 分裂. 因而可用 Raynaud 的诸多下降有效性判据之一 [23, XI~3.1(6)] 得到结论.

\medskip
\noindent
\SGARefTarget{sga7.sga7i.eix.n7-2-1-4}{7.2.1.4}. 这就完成了 (b) 的证明, 只剩下全忠实性断言 (a). 留给读者验证, 这一断言实质上已包含在 (b) 的证明方法中; 或者读者也可自行寻找独立的直接证明.

\medskip
\noindent
\SGARefTarget{sga7.sga7i.eix.n7-2-1-5}{7.2.1.5}. 注意, 若假设 $A^o$ 在 $S$ 上拟射影, 或等价地, $B$ 在 $S$ 上射影, 则无需引用 Raynaud 精细的下降结果. 特别地, 若从 $S$ 上的形式群概形 $\widehat A$ 出发, 其中 $A$ 如 \SGARefLink{sga7.sga7i.eix.n7-1-5}{7.1.5}, 则此条件成立. 这是我们使用 \SGARefLink{sga7.sga7i.eix.n7-2-1}{7.2.1} 的唯一情形.

% END INCLUDED SGA7I_English_current_buildable_source_20260823/components/150_expose_IX_lemma_7_2_1.tex

% BEGIN INCLUDED SGA7I_English_current_buildable_source_20260823/components/151_expose_IX_sections_7_2_2_through_7_2_6.tex
% Expose IX, 7.2.2--7.2.6.
% French transcription lines 1940--1967; authority scan physical pp. 418--419
% (scan indices 417--418; source folios 406--407).  Both pages were checked
% directly in four overlapping 1100-dpi bands; no escalation was needed.

\medskip
\noindent
\SGARefTarget{sga7.sga7i.eix.n7-2-2}{7.2.2}. 将 \SGARefLink{sga7.sga7i.eix.n7-2-1}{7.2.1} 用于 \SGARefLink{sga7.sga7i.eix.n7-1}{7.1} 中考虑的形式群概形 $\widehat A$, 取该处处理的有利情形, 即形式 阿贝尔概形 $\widehat B$ 可代数化 (参见 \SGARefLink{sga7.sga7i.eix.n7-1-5}{7.1.5}), 并额外假设 $S_o$ 上的环面 $T_o$ 等平凡. 于是得到 $S$ 上的群概形 $A^\natural$, 称为 $A$ 在 $S$ 上的 \underline{Raynaud 群}; 它精确到同构, 由 $\operatorname{Spf}(V)=\widehat S$ 上形式群概形的一个典范同构刻画:
\begin{equation}\tag{\SGARefTarget{sga7.sga7i.eix.n7-2-3}{7.2.3}}
 (A^\natural)^\wedge\simeq\widehat A.
\end{equation}
由 \SGARefLink{sga7.sga7i.eix.n7-2-1}{7.2.1},
\begin{equation}\tag{\SGARefTarget{sga7.sga7i.eix.n7-2-4}{7.2.4}}
 \Phi=A^\natural/A^{\natural o}
\end{equation}
为 $S$ 上的有限平展群, 而单位元连通分支
\[
 A^{\natural o}=A^{o\natural}
\]
具有扩张结构
\begin{equation}\tag{\SGARefTarget{sga7.sga7i.eix.n7-2-5}{7.2.5}}
 0\longrightarrow T\longrightarrow A^{\natural o}
 \longrightarrow B\longrightarrow0,
\end{equation}
其中 $B$ 为 $S$ 上的 阿贝尔概形, 其形式完备化就是 \SGARefLink{sga7.sga7i.eix.n7-1}{7.1} 中构造的形式 阿贝尔概形 $\widehat B$, 而 $T$ 则精确到唯一同构, 由它是 $S$ 上提升 $T_o$ 的等平凡环面这一性质刻画. 限制到各 $S_n$ 时, 扩张结构 (\SGARefLink{sga7.sga7i.eix.n7-2-5}{7.2.5}) 给出 (\SGARefLink{sga7.sga7i.eix.n7-1-1}{7.1.1}) 在 $A_n^{\natural o}=A_n^o$ 上定义的扩张结构.

\medskip
\noindent
\SGARefTarget{sga7.sga7i.eix.n7-2-6}{7.2.6}. 由 \SGARefLink{sga7.sga7i.eix.n7-2-1}{7.2.1} 也立即得到, Raynaud 群 $A^\natural$ 以及正合列 (\SGARefLink{sga7.sga7i.eix.n7-2-5}{7.2.5}), 均函子性地依赖于 $A$ (满足当前考虑的条件), 且其构成与理想进拓扑环同态 $V\longrightarrow V'$ 所对应的任意基变换 $S'\longrightarrow S$ 交换.

% END INCLUDED SGA7I_English_current_buildable_source_20260823/components/151_expose_IX_sections_7_2_2_through_7_2_6.tex

% BEGIN INCLUDED SGA7I_English_current_buildable_source_20260823/components/152_expose_IX_section_7_3.tex
% Expose IX, section 7.3.
% French transcription lines 1969--2015; authority scan physical pp. 419--421
% (scan indices 418--420; source folios 407--409). All three pages were checked
% directly in four overlapping 1100-dpi bands; no escalation was needed.

\medskip
\noindent
\SGARefTarget{sga7.sga7i.eix.n7-3}{7.3}. 现在对固定素数 $\ell$, 考虑 pro-$\ell$ 群 $T_{\ell}(A^{o})$. 将 $S$ 上的有限概形与相应的形式概形等同, 并如常定义「\underline{不变部分}」$T_{\ell}(A^{o})^{f}$, 便得到如 (\SGARefLink{sga7.sga7i.eix.n5-1-5}{5.1.5}) 中的同构:

\begin{equation}\tag{\SGARefTarget{sga7.sga7i.eix.n7-3-1}{7.3.1}}
T_{\ell}(A^{o})^{f} \simeq T_{\ell}(\widehat{A}^{o}) \quad ,
\end{equation}
再如 (\SGARefLink{sga7.sga7i.eix.n5-1-6}{5.1.6}) 定义「\underline{环面部分}」,

\begin{equation}\tag{\SGARefTarget{sga7.sga7i.eix.n7-3-2}{7.3.2}}
T_{\ell}(A^{o})^{t} = T_{\ell}(\widehat{T}) \simeq T_{\ell}(T)
 \subset T_{\ell}(A^{o})^{f} \simeq T_{\ell}(\widehat{A}^{o}) \quad ,
\end{equation}
立即对「\underline{阿贝尔部分}」得到同构

\begin{equation}\tag{\SGARefTarget{sga7.sga7i.eix.n7-3-3}{7.3.3}}
T_{\ell}(A^{o})^{ab} \stackrel{\mathrm{def}}{=}
T_{\ell}(A^{o})^{f}/T_{\ell}(A^{o})^{t}
 \simeq T_{\ell}(\widehat{B}) \quad .
\end{equation}
这里尚未使用 $\widehat B$ 可代数化这一假设; 正是该假设使我们能够构造 阿贝尔概形 $B$ 及 $B$ 被 $T$ 的 Raynaud 扩张. 当然存在典范同构 (即 (\SGARefLink{sga7.sga7i.eix.n7-3-1}{7.3.1}) 应用于 阿贝尔概形 $B$ 的特例!) $T_{\ell}(\widehat B)\simeq T_{\ell}(B)$, 因此同构 (\SGARefLink{sga7.sga7i.eix.n7-3-3}{7.3.3}) 还给出

\begin{equation}\tag{\SGARefTarget{sga7.sga7i.eix.n7-3-4}{7.3.4}}
T_{\ell}(A^{o})^{ab} \simeq T_{\ell}(B) \quad .
\end{equation}
顺便注意, (\SGARefLink{sga7.sga7i.eix.n7-2-3}{7.2.3}) 与 (\SGARefLink{sga7.sga7i.eix.n7-3-1}{7.3.1}) 蕴涵典范同构

\begin{equation}\tag{\SGARefTarget{sga7.sga7i.eix.n7-3-5}{7.3.5}}
T_{\ell}(A^{o})^{f} \simeq T_{\ell}(A^{\natural o}) \quad ,
\end{equation}
因为右端已经是投射有限群, 即与其不变部分相同. 典范正合列

\begin{equation}\tag{\SGARefTarget{sga7.sga7i.eix.n7-3-6}{7.3.6}}
0 \longrightarrow T_{\ell}(A^{o})^{t}
 \longrightarrow T_{\ell}(A^{o})^{f}
 \longrightarrow T_{\ell}(B) \longrightarrow 0 \quad ,
\end{equation}
在任何情形下 (不要求 $\widehat B$ 可代数化), 都可解释为将函子 $T_{\ell}$ 用于正合列

\begin{equation*}
0 \longrightarrow \widehat{T} \longrightarrow \widehat{A}^{o}
 \longrightarrow \widehat{B} \longrightarrow 0 \quad ,
\end{equation*}
所得的正合列; 它也可看作将函子 $T_{\ell}$ 用于通常的群概形的典范正合列

\begin{equation*}
0 \longrightarrow T \longrightarrow A^{\natural o}
 \longrightarrow B \longrightarrow 0 \qquad .
\end{equation*}

% END INCLUDED SGA7I_English_current_buildable_source_20260823/components/152_expose_IX_section_7_3.tex

% BEGIN INCLUDED SGA7I_English_current_buildable_source_20260823/components/153_expose_IX_section_7_4.tex
% Expose IX, section 7.4.
% French transcription lines 2017--2080; authority scan physical pp. 421--424
% (scan indices 420--423; source folios 409--412). All four pages were
% checked directly in four overlapping 1100-dpi bands; no escalation was needed.

\medskip
\noindent
\SGARefTarget{sga7.sga7i.eix.n7-4}{7.4}. 现在设还给定 $S$ 上另一个群概形 $A'$, 满足与 $A$ 相同的假设, 并假设 $A$ 和 $A'$ 均具有连通纤维. 设 $W$ 为 $(A,A')$ 被 $\mathbb G_{mS}$ 的一个双扩张. 于是得到配对 (VIII \SGARefLink{sga7.sga7i.eviii.n2-2}{2.2})

\begin{equation}\tag{\SGARefTarget{sga7.sga7i.eix.n7-4-1}{7.4.1}}
T_{\ell}(A) \times T_{\ell}(A') \longrightarrow
T_{\ell}(\mathbb G_{mS}) = \mathbb Z_{\ell}(1)
\end{equation}
再限制到不变部分, 得到配对

\begin{equation}\tag{\SGARefTarget{sga7.sga7i.eix.n7-4-2}{7.4.2}}
T_{\ell}(A)^{f} \times T_{\ell}(A')^{f}
 = T_{\ell}(\widehat A) \times T_{\ell}(\widehat {A'})
 \longrightarrow \mathbb Z_{\ell}(1) \qquad .
\end{equation}
这也可解释为将 $W$ 完备化得到的 $(\widehat A,\widehat {A'})$ 被 $\widehat{\mathbb G}_{mS}$ 的双扩张 $\widehat W$ 所定义的配对 (为简便起见, 这里将形式群概形的双扩张定义为 $(A_n,A'_n)$ 被 $\mathbb G_{mS_n}$ 的双扩张 $W_n$ 的相容族). 已知 (VIII \SGARefLink{sga7.sga7i.eviii.n3-5}{3.5}), $W_n$ 来自 $(B_n,B'_n)$ 的一个双扩张 $W_n^{ab}$, 其中 $B'_n$ 为 $A'_n$ 的 阿贝尔部分, 因而 (\SGARefLink{sga7.sga7i.eix.n7-4-2}{7.4.2}) 通过下述配对分解

\begin{equation}\tag{\SGARefTarget{sga7.sga7i.eix.n7-4-3}{7.4.3}}
T_{\ell}(A)^{ab} \times T_{\ell}(A')^{ab}
 = T_{\ell}(\widehat B) \times T_{\ell}(\widehat {B'})
 \longrightarrow \mathbb Z_{\ell}(1) \qquad .
\end{equation}
这也是双扩张

\begin{equation*}
\widehat W^{ab}=(W_n^{ab})
\end{equation*}
所定义的配对, 该双扩张是 $(\widehat B,\widehat {B'})$ 被 $\widehat{\mathbb G}_{mS}$ 的双扩张. 当 $\widehat B$ 和 $\widehat {B'}$ 可代数化, 即存在 $S$ 上的 阿贝尔概形 $B,B'$ 时, EGA IV 5.1.4 表明 $\widehat W^{ab}$ 来自一个精确到唯一同构被定义的双扩张

\begin{equation*}
W^{ab}\quad\text{为}\quad(B,B')\quad\text{被}\quad\mathbb G_{mS},
\end{equation*}
且相应的配对

\begin{equation}\tag{\SGARefTarget{sga7.sga7i.eix.n7-4-4}{7.4.4}}
T_{\ell}(B) \times T_{\ell}(B')
 \longrightarrow T_{\ell}(\mathbb G_{mS})
\end{equation}
典范同构于 (\SGARefLink{sga7.sga7i.eix.n7-4-3}{7.4.3}).

以 $D(B')$ 表示 $B'$ 的对偶 阿贝尔概形, 则 $W^{ab}$ 定义一个同态

\begin{equation}\tag{\SGARefTarget{sga7.sga7i.eix.n7-4-5}{7.4.5}}
u:B\longrightarrow D(B') \qquad ,
\end{equation}
其性质反映了配对 (\SGARefLink{sga7.sga7i.eix.n7-4-4}{7.4.4}), 即 (\SGARefLink{sga7.sga7i.eix.n7-4-3}{7.4.3}) 的性质. 例如, 配对 (\SGARefLink{sga7.sga7i.eix.n7-4-3}{7.4.3}) 在点 $t\in S$ 处为完美对偶, 当且仅当 $u_t:B_t\longrightarrow D(B')_t$ 是核的秩素于 $\ell$ 的同源. 这蕴涵 $u$ 本身是同源 (因为 $S$ 局部, 因而连通), 且 $\ker u$ 在 $S$ 上平坦且有限, 所以 $\ker u$ 在每个点的秩均相同. 因而刚才考虑的条件与所选点 $t$ 无关.

\medskip
\noindent
\SGARefTarget{sga7.sga7i.eix.n7-4-6}{7.4.6}. 特别地, 这完成了 \SGARefLink{sga7.sga7i.eix.n5-4}{5.4} 的证明. 事实上, 在该定理的条件下 (可设 $S$ 完备), 将上述讨论用于 $A^o$ 与 ${A'}^o$, 可知对每个 $\ell$, 配对 (\SGARefLink{sga7.sga7i.eix.n7-4-3}{7.4.3}) 在一般点 $\eta$ 处为完美对偶, 因而在点 $s$ 处也如此.

\medskip
\noindent
\SGARefTarget{sga7.sga7i.eix.n7-4-7}{7.4.7}. 上述讨论也使正交性定理 \SGARefLink{sga7.sga7i.eix.n5-2}{5.2} 清楚明白, 无需借助引理 \SGARefLink{sga7.sga7i.eix.n5-2-4}{5.2.4}. 事实上, 在 \SGARefLink{sga7.sga7i.eix.n7-4}{7.4} 中已经顺便证明, 配对 $(\SGARefLink{sga7.sga7i.eix.n5-2-3-1}{5.2.3.1})=(\SGARefLink{sga7.sga7i.eix.n7-4-2}{7.4.2})$ 在两边的环面部分上均为零, 因而通过配对 (\SGARefLink{sga7.sga7i.eix.n7-4-3}{7.4.3}) 分解. 已知该配对在所考虑投射群的特殊纤维上非退化, 只需证明它在一般纤维上也非退化. 这来自以下事实: 若 (\SGARefLink{sga7.sga7i.eix.n7-4-5}{7.4.5}) 在特殊纤维上诱导同源, 则它本身就是同源.

\medskip
\noindent
\SGARefTarget{sga7.sga7i.eix.n7-4-8}{7.4.8}. 回到 \SGARefLink{sga7.sga7i.eix.n7-4}{7.4} 的一般条件, 并设给定开子集 $U\subset S-S_o$, 使 $A_U$ 与 $A'_U$ 为通过双扩张 $W_U$ 彼此对偶的 \underline{阿贝尔概形}. 将 (\SGARefLink{sga7.sga7i.eix.n7-3-6}{7.3.6}) 以及 $A'$ 的类似正合列限制到 $U$, 得到 Barsotti--Tate pro-$\ell$ 群 $T_\ell(A_U)$ 与 $T_\ell(A'_U)$ 上的典范滤过:

\begin{equation}\tag{\SGARefTarget{sga7.sga7i.eix.n7-4-9}{7.4.9}}
\left\{
\begin{aligned}
T_\ell(A_U)&\supset T_\ell(A_U)^f\supset T_\ell(A_U)^t\supset0,\\[4ex]
T_\ell(A'_U)&\supset T_\ell(A'_U)^f\supset T_\ell(A'_U)^t\supset0
\qquad .
\end{aligned}
\right.
\end{equation}
断言\underline{这两个滤过关于配对 (\SGARefLink{sga7.sga7i.eix.n7-4-1}{7.4.1}) 彼此对偶}, 此配对在这里是完美对偶 (VIII \SGARefLink{sga7.sga7i.eviii.n3-2}{3.2}); 即将 $A_K,A'_K$ 换成 $A_U,A'_U$ 后, 关系式 (\SGARefLink{sga7.sga7i.eix.n5-2-2}{5.2.2}) 及其对称的关系仍然成立. 上面已看到, 此正交性断言还等价于说, $(B,B')$ 被 $\mathbb G_{mS}$ 的双扩张 $W^{ab}$ 在 $U$ 上为完美对偶 (或者当 $U$ 在 $S$ 中稠密时, 等价地在 $S$ 上为完美对偶). 只需在 $U$ 的每个点 $\eta$ 处验证这一性质; 于是通常的论证归结到已知的情形, 即 $S$ 为一般点为 $\eta$ 的性状.

正交性关系还允许将公式 (\SGARefLink{sga7.sga7i.eix.n5-5-7}{5.5.7})、(\SGARefLink{sga7.sga7i.eix.n5-5-8}{5.5.8}) 及 (\SGARefLink{sga7.sga7i.eix.n5-5-9}{5.5.9}) 推广到当前情形, 只需以限制到开子集 $U$ 代替限制到一般点 $\eta$. 因而 \SGARefLink{sga7.sga7i.eix.n5-6}{5.6} 同样推广到当前情形. 这里不在这一背景下重写上述公式与陈述.

% END INCLUDED SGA7I_English_current_buildable_source_20260823/components/153_expose_IX_section_7_4.tex

% BEGIN INCLUDED SGA7I_English_current_buildable_source_20260823/components/154_expose_IX_proposition_7_5_and_section_7_6.tex
% Expose IX, Proposition 7.5 and section 7.6.
% French transcription lines 2082--2103; authority scan physical pp. 424--425
% (scan indices 423--424; source folios 412--413). Both pages were checked
% directly in four overlapping 1100-dpi bands; no escalation was needed.

从 N\'{e}ron 模型的观点, 汇总所得到的关于 阿贝尔部分及其配对的结果.

\medskip
\noindent
\underline{命题} \SGARefTarget{sga7.sga7i.eix.n7-5}{7.5}. \underline{设 $S$ 为函数域为 $K$ 的完备性状. 对 $K$ 上每个在 $S$ 上具有半稳定约化的 阿贝尔概形 $A_K$, 将 \SGARefLink{sga7.sga7i.eix.n7-1}{7.1} 的过程应用于 $A_K$ 的 N\'{e}ron 模型 $A$ 的单位元连通分支 $A^o$, 就给它对应一个 阿贝尔部分 $B=A^{ab}$, 它是 $S$ 上函子性地依赖于 $A_K$ 的 阿贝尔概形. 对每个具有半稳定约化的另一 阿贝尔概形 $A'_K$, 可给 $(A_K,A'_K)$ 上每个除子对应 $W_K$ 对应一个 $(A^{ab},{A'}^{ab})$ 上唯一的除子对应 $W^{ab}$, 使得对每个整数 $n\geq0$, $W_n^{ab}$ 沿 $A_n^o\longrightarrow B_n=A_n^{oab}$ 与 ${A'}_n^o\longrightarrow B'_n={A'}_n^{oab}$ 的拉回为由 $W$ 约化所得的 $(A_n^o,{A'}_n^o)$ 的双扩张 $W_n$, 其中 $W$ 是 $W_K$ 到 $(A^o,{A'}^o)$ 的典范延拓 (\SGARefLink{sga7.sga7i.eix.n1-4}{1.4}).

$W^{ab}$ 的构成与沿 阿贝尔概形态射 $\overline A_K\longrightarrow A_K$ 和 $\overline {A'}_K\longrightarrow A'_K$ 取对应的拉回相容. 若 $W_K$ 为完美对偶, 则 $W^{ab}$ 也如此. 若 $W_K$ 为与 $A_K$ 的某个极化关联的 $(A_K,A_K)$ 上的除子对应, 则 $W^{ab}$ 与 $A^{ab}$ 的某个极化关联.}

$W^{ab}$ 的构成与拉回的相容性立即来自上面对 $W^{ab}$ 的刻画. 关于极化的额外断言已在 \SGARefLink{sga7.sga7i.eix.n7-1-5}{7.1.5} 中与 $\widehat B$ 的可代数化性一起证明.

\medskip
\noindent
\SGARefTarget{sga7.sga7i.eix.n7-6}{7.6}. 继续假设 $S$ 为完备性状, $A$ 为 $K$ 上 阿贝尔概形 $A_K$ 的 N\'{e}ron 模型. 置

\begin{equation}\tag{\SGARefTarget{sga7.sga7i.eix.n7-6-1}{7.6.1}}
A_K^{\natural}=(A^{\natural})_K \quad ,
\end{equation}
其中 $A^{\natural}$ 为 Raynaud 群概形 (\SGARefLink{sga7.sga7i.eix.n7-2-2}{7.2.2}): 它是有限平展群 $\Phi$ 被 $A^{\natural o}$ 的扩张, 而 $A^{\natural o}$ 是 阿贝尔概形 $B$ 被环面 $T$ 的扩张. 因而 $A_K^{\natural}$ 本身是 $\Phi_K$ 被 $(A_K^{\natural})^o=(A^{\natural o})_K=(A^{o\natural})_K$ 的扩张, 后者是 阿贝尔概形 $B_K$ 被环面 $T_K$ 的扩张. 对每个素数 $\ell$, 由 (\SGARefLink{sga7.sga7i.eix.n7-3-5}{7.3.5}) 及 (\SGARefLink{sga7.sga7i.eix.n7-3-2}{7.3.2}) 得到

\begin{equation}\tag{\SGARefTarget{sga7.sga7i.eix.n7-6-2}{7.6.2}}
T_{\ell}(A_K^{\natural}) \simeq T_{\ell}(A_K^{\natural o})
 \simeq T_{\ell}(A_K)^f \; , \qquad
T_{\ell}(T_K) \simeq T_{\ell}(A_K)^t \quad .
\end{equation}

\medskip
\noindent
\SGARefTarget{sga7.sga7i.eix.n7-6-3}{7.6.3}. 由 \SGARefLink{sga7.sga7i.eix.n3-3}{3.3}, $A_K^{\natural o}$ 的构成与完备性状态射 $S'\longrightarrow S$ 所对应的任意基变换交换.

% END INCLUDED SGA7I_English_current_buildable_source_20260823/components/154_expose_IX_proposition_7_5_and_section_7_6.tex

% BEGIN INCLUDED SGA7I_English_current_buildable_source_20260823/components/155_expose_IX_section_8_1.tex
% Expose IX, section 8.1.
% French transcription lines 2109--2154; authority scan physical pp. 426--428
% (scan indices 425--427; source folios 414--416). All three pages were
% checked directly in four overlapping 1100-dpi bands; no escalation was
% needed.

\medskip
\noindent
8. \underline{非半稳定情形中的 Raynaud 扩张 $A_K^{\natural o}$, 以及归纳群 $A_K^{\flat}$}

\medskip
\noindent
\SGARefTarget{sga7.sga7i.eix.n8-1}{8.1}. 继续假设性状 $S$ 完备, 但不再要求 $A_K$ 必有半稳定约化. 如常令 $K'$ 为 $K$ 的有限 Galois 扩张, 使 $A_{K'}$ 关于 $S$ 在 $K'$ 中的正规化 $S'$ 有半稳定约化. 考虑与 $A_{K'}$ 关联的 Raynaud 扩张 $A_{K'}^{\natural}$ (\SGARefLink{sga7.sga7i.eix.n7-6-1}{7.6.1}). 显然, $K'/K$ 的 Galois 群通过结构搬移作用在此群概形上; 因而由 Galois 下降, 它给出 $K$ 上的一个群概形, 记为 $A_{K,K'}^{\natural}$, 它依赖于扩张 $K'$ 的选取. 然而由 \SGARefLink{sga7.sga7i.eix.n7-6-3}{7.6.3}, 其单位元连通分支精确到典范同构与此选取无关, 简记为 $A_K^{\natural o}$. 下降表明 $A_K^{\natural o}$ 是一个扩张

\begin{equation}\tag{\SGARefTarget{sga7.sga7i.eix.n8-1-1}{8.1.1}}
0 \longrightarrow T_K \longrightarrow A_K^{\natural o}
 \longrightarrow B_K \longrightarrow 0 \quad ,
\end{equation}
其中 $T_K$ 为环面, $B_K$ 为 $K$ 上的 阿贝尔概形. 类似地, 根据 $T_{\ell}(A_K)$ 的本质不变部分和本质环面部分的定义 (\SGARefLink{sga7.sga7i.eix.n4-1}{4.1} 以及 6.5), 同构 (\SGARefLink{sga7.sga7i.eix.n7-6-2}{7.6.2}) 给出典范同构

\begin{equation}\tag{\SGARefTarget{sga7.sga7i.eix.n8-1-2}{8.1.2}}
T_{\ell}(A_K^{\natural}) \simeq T_{\ell}(A_K)^{ef} \quad ,
\qquad T_{\ell}(T_K) \simeq T_{\ell}(A_K)^{et} \quad ,
\end{equation}
从而取商得到

\begin{equation}\tag{\SGARefTarget{sga7.sga7i.eix.n8-1-3}{8.1.3}}
T_{\ell}(B_K) \simeq
T_{\ell}(A_K)^{ef}/T_{\ell}(A_K)^{et} \quad .
\end{equation}

这样就确定了滤过

\begin{equation}\tag{\SGARefTarget{sga7.sga7i.eix.n8-1-4}{8.1.4}}
T_\ell(A_K) \supset T_\ell(A_K)^{ef}
 \supset T_\ell(A_K)^{et} \supset 0
\end{equation}
中的两个商. 因而群概形扩张 (\SGARefLink{sga7.sga7i.eix.n8-1-1}{8.1.1}) 描述了此滤过中的两个商. 对第三个商 $T_\ell(A_K)/T_\ell(A_K)^{ef}$, (6.5.2) 给出典范同构

\begin{equation}\tag{\SGARefTarget{sga7.sga7i.eix.n8-1-5}{8.1.5}}
T_\ell(A_K)/T_\ell(A_K)^{ef}
 \simeq D(T_\ell(A'_K)^{et}) \quad ,
\end{equation}
其中 $D=\underline{\operatorname{Hom}}(-,T_\ell(\mathbb{G}_m))$ 表示 Cartier 对偶 (在 Barsotti--Tate 投射群的意义下), $A'_K$ 表示 $A_K$ 的对偶 阿贝尔概形. 再考虑与 (\SGARefLink{sga7.sga7i.eix.n8-1-1}{8.1.1}) 同型的正合列

\begin{equation}\tag{\SGARefTarget{sga7.sga7i.eix.n8-1-6}{8.1.6}}
0 \longrightarrow T'_K \longrightarrow {A'_K}^{\natural o}
 \longrightarrow B'_K \longrightarrow 0
\end{equation}
它与 $K$ 上的 阿贝尔概形 $A'_K$ 关联, 并考虑相应的 (\SGARefLink{sga7.sga7i.eix.n8-1-2}{8.1.2}) 型同构. 于是公式 (\SGARefLink{sga7.sga7i.eix.n8-1-5}{8.1.5}) 给出

\begin{equation}\tag{\SGARefTarget{sga7.sga7i.eix.n8-1-7}{8.1.7}}
T_\ell(A_K)/T_\ell(A_K)^{ef}
 \simeq D(T_\ell(T'_K))
 \simeq \underline{M}_K\otimes_{\mathbb{Z}}\mathbb{Z}_\ell \quad ,
\end{equation}
其中推广 (\SGARefLink{sga7.sga7i.eix.n5-5-2}{5.5.2}) 中引入的记号, 置

\begin{equation}\tag{\SGARefTarget{sga7.sga7i.eix.n8-1-8}{8.1.8}}
\underline{M}_K=D(T'_K) \quad , \qquad
\underline{M}'_K=D(T_K) \quad ,
\end{equation}
其中 $D$ 现在表示通常群概形意义下的 Cartier 对偶 $\underline{\operatorname{Hom}}(-,\mathbb{G}_{mK})$. 因而 $\underline{M}_K$ 与 $\underline{M}'_K$ 是 $K$ 上的扭常值群, 它们在 $K'$ 上成为常值群, 同构于形如 $\mathbb{Z}^{r}$ 的群, 其中 $r=\dim T=\dim T'$.

% END INCLUDED SGA7I_English_current_buildable_source_20260823/components/155_expose_IX_section_8_1.tex

% BEGIN INCLUDED SGA7I_English_current_buildable_source_20260823/components/156_expose_IX_sections_8_2_and_8_3.tex
% Expose IX, sections 8.2--8.3.
% French transcription lines 2159--2188; authority scan physical pp. 428--429
% (scan indices 427--428; source folios 416--417). Both pages were checked
% directly in four overlapping 1100-dpi bands; physical p. 430 was checked
% for the exact section-9 boundary. The authority's evident A_K -> A'_K
% polarization is restored where the working transcription misplaces the
% prime as A_{K'}; its evident "extensions ... de K'" slip is rendered
% according to the mathematically required extensions of K.

\medskip
\noindent
\SGARefTarget{sga7.sga7i.eix.n8-2}{8.2}. 假设如 \SGARefLink{sga7.sga7i.eix.n8-1}{8.1}. 显然, 扩张 (\SGARefLink{sga7.sga7i.eix.n8-1-1}{8.1.1}) 函子性地依赖于 阿贝尔概形 $A_K$, 且其构成与任意性状态射 $\widetilde S\longrightarrow S$ 交换. 选取 $A_K$ 的一个极化, 从而得到同源 $A_K\longrightarrow A'_K$. 它诱导从 (\SGARefLink{sga7.sga7i.eix.n8-1-1}{8.1.1}) 到 (\SGARefLink{sga7.sga7i.eix.n8-1-6}{8.1.6}) 的同态; 观察 $T_\ell$ 就立即看出, 此同态逐项为同源. 类似地, 利用 \SGARefLink{sga7.sga7i.eix.n7-5}{7.5} 与下降, 可知 $K$ 上 阿贝尔概形 $A_K$ 与 $\overline A_K$ 被 $\mathbb{G}_{mK}$ 的每个双扩张, 都定义 $(B_K,\overline B_K)$ 被 $\mathbb{G}_{mK}$ 的一个双扩张, 其中 $B_K$ 与 $\overline B_K$ 分别表示对应于 $A_K$ 和 $\overline A_K$ 的 Raynaud 扩张的 阿贝尔商. 当 $W$ 为完美对偶时, $(B_K,\overline B_K)$ 的相应双扩张也如此. 特别地, 沿用 \SGARefLink{sga7.sga7i.eix.n7-6}{7.6} 的记号, 阿贝尔概形 $B_K$ 和 $B'_K$ 典范地互为对偶.

当 $A_K$ 与 $\overline A_K$ 变化时, $(A_K,\overline A_K)$ 上的双扩张与 $(B_K,\overline B_K)$ 上的双扩张之间的对应, 具有显然的函子性. 它们由 \SGARefLink{sga7.sga7i.eix.n7-5}{7.5} 中陈述的类似性质经下降得到; 其表述留给读者. 也将上述构造与任意性状变换 $\widetilde S\longrightarrow S$ 的相容性留给读者.

\medskip
\noindent
\SGARefTarget{sga7.sga7i.eix.n8-3}{8.3}. 回到群 $A_{K,K'}^{\natural}$, 研究它随 $K'$ 的变化. 显然, $K$ 的两个同构扩张 $K'$ 与 $K'_1$ 给出典范同构的群, 因而只需考虑 $K$ 的可分闭包 $\overline K$ 中的子扩张 $K'$. 若 $K'\hookrightarrow K''$ 为 $\overline K$ 中两个这样的有限 Galois 子扩张, 则由 \SGARefLink{sga7.sga7i.eix.n3-3}{3.3} 立即得到典范\underline{开浸入}

\begin{equation}\tag{\SGARefTarget{sga7.sga7i.eix.n8-3-1}{8.3.1}}
A^{\natural}_{K,K'} \hookrightarrow A^{\natural}_{K,K''}
 \qquad ,
\end{equation}
它将第一个群等同于第二个群的一个开子群. 因而当 $K'$ 变化时, 各 $A^{\natural}_{K,K'}$ 组成 $K$ 上有限型光滑群的归纳系统, 其过渡态射均为开浸入. 若置

\begin{equation}\tag{\SGARefTarget{sga7.sga7i.eix.n8-3-2}{8.3.2}}
\Phi_{K,K'}=A^{\natural}_{K,K'}
 /{A^{\natural}_{K,K'}}^{o} \qquad ,
\end{equation}
则各 $\Phi_{K,K'}$ 自身组成 $K$ 上有限平展群的归纳系统, 下文将对此作更详细研究.

置

\begin{equation}\tag{\SGARefTarget{sga7.sga7i.eix.n8-3-3}{8.3.3}}
A^{\flat}_{K}=\varinjlim_{K'} A^{\natural}_{K,K'} \qquad ,
\end{equation}
其中极限取在 $K$ 上的 fppf 层范畴中. 此外, 由于此归纳系统的过渡态射均为开浸入, 显然此极限可由 $K$ 上的一个群概形表示 (一般并不拟紧), 自然仍以同一个符号 $A^{\flat}_{K}$ 表示它. 对此群概形有

\begin{equation}\tag{\SGARefTarget{sga7.sga7i.eix.n8-3-4}{8.3.4}}
{A_{K}^{\flat}}^{o}=A_{K}^{\natural o} \quad , \qquad
A^{\flat}_{K}/{A_{K}^{\flat}}^{o}
 \simeq \varinjlim_{K'}\Phi_{K,K'} \qquad .
\end{equation}
后一个群将在 \SGARefLink{sga7.sga7i.eix.n11-10}{11.10} 中确定: 沿用 (\SGARefLink{sga7.sga7i.eix.n8-1-8}{8.1.8}) 的记号, 它典范同构于 $\underline{M}_{K}\otimes\mathbb{Q}/\mathbb{Z}$.

% END INCLUDED SGA7I_English_current_buildable_source_20260823/components/156_expose_IX_sections_8_2_and_8_3.tex

% BEGIN INCLUDED SGA7I_English_current_buildable_source_20260823/components/157_expose_IX_section_9_1.tex
% Expose IX, section 9 heading and 9.1.
% French transcription lines 2194--2214; authority scan physical pp. 430--431
% (scan indices 429--430; source folios 418--419). Both pages were checked
% directly in four overlapping 1100-dpi bands; no escalation was needed.

\medskip
\noindent
9. \underline{单值配对的定义}
$\underline{M}_{\ell}\otimes\underline{M}'_{\ell}
 \longrightarrow \mathbb{Z}_{\ell S}$.

\medskip
\noindent
\SGARefTarget{sga7.sga7i.eix.n9-1}{9.1}. 本节假设 $S$ 为 Hensel 性状, 且 $A_K$ 在 $S$ 上有半稳定约化, 或等价地, $A'_K$ 在 $S$ 上有半稳定约化 (\SGARefLink{sga7.sga7i.eix.n3-5-1}{3.5.1}). 于是 $S$ 上有两个扭常值 $\mathbb{Z}$-模 $\underline{M}$ 与 $\underline{M}'$ (\SGARefLink{sga7.sga7i.eix.n5-5-2}{5.5.2}). 在 $S$ 的适当有限平展覆盖上 (因而若 $S$ 严格局部, 即 $k$ 可分闭, 则在 $S$ 本身上), 它们分别分裂为取值于 $M$ 与 $M'$ 的常值群, 非典范同构于 $\mathbb{Z}^{r}$, 其中 $r$ 为 $A_o$ 和 $A'_o$ 的极大环面的共同维数. 对每个素数 $\ell$, 将考虑相应的 $\ell$ 进层

\begin{equation}\tag{\SGARefTarget{sga7.sga7i.eix.n9-1-1}{9.1.1}}
\underline{M}_{\ell}=\underline{M}\otimes_{\mathbb{Z}}\mathbb{Z}_{\ell}
 \quad , \qquad
\underline{M}'_{\ell}=\underline{M}'\otimes_{\mathbb{Z}}\mathbb{Z}_{\ell}
 \qquad ,
\end{equation}
其几何意义由公式 (\SGARefLink{sga7.sga7i.eix.n5-5-7}{5.5.7}) 及 (\SGARefLink{sga7.sga7i.eix.n5-5-8}{5.5.8}) 明确给出. 本节将定义一个典范配对, 称为\underline{单值配对},

\begin{equation}\tag{\SGARefTarget{sga7.sga7i.eix.n9-1-2}{9.1.2}}
u_{\ell}:\underline{M}_{\ell}\otimes\underline{M}'_{\ell}
 \longrightarrow\mathbb{Z}_{\ell} \qquad ,
\end{equation}
它将由 Barsotti--Tate pro-$\ell$ 群 $T_{\ell}(A^{o})$ 与 $T_{\ell}({A'}^{o})$ 定义, 甚至可由第一个群的一般纤维 $T_{\ell}(A_K)$ 定义 (当 $\ell=p$ 时, 要求 Tate 定理 [33, 定理~4] 可用, 即就目前而言须处于混合特征情形). 当 $\ell\neq p$ 时, 这等价于说配对 (\SGARefLink{sga7.sga7i.eix.n9-1-2}{9.1.2}) 可由 (\SGARefLink{sga7.sga7i.eix.n2-5-3}{2.5.3}) 中考虑的 Galois $\pi$-模

\begin{equation}\tag{\SGARefTarget{sga7.sga7i.eix.n9-1-3}{9.1.3}}
U=T_{\ell}(A_K)(\overline K)=T_{\ell}(A_{\overline K})
\end{equation}
来描述. 更精确地, 我们将看到, 知道这一配对, 实质上等价于知道单值群 $I$ 在 $U$ 上的作用.

% END INCLUDED SGA7I_English_current_buildable_source_20260823/components/157_expose_IX_section_9_1.tex

% BEGIN INCLUDED SGA7I_English_current_buildable_source_20260823/components/158_expose_IX_section_9_2.tex
% Expose IX, section 9.2.
% French transcription lines 2216--2276; authority scan physical pp. 431--433
% (scan indices 430--432; source folios 419--421). All three pages were
% checked directly in four overlapping 1100-dpi bands; physical p. 434 was
% checked for the exact 9.3 boundary. The authority controls the evident
% transcription repairs x (not u) in the generator phrase, V^\perp (not
% V^1) in (9.2.4), T'_o in the second line of (9.2.7), and balanced
% parentheses in the displayed formulas.

\medskip
\noindent
\SGARefTarget{sga7.sga7i.eix.n9-2}{9.2}. \underline{$\ell\neq p$ 的情形}. 从较简单的 $\ell\neq p$ 情形开始, 沿用 (\SGARefLink{sga7.sga7i.eix.n2-5-3}{2.5.3}) 中引入的记号. 注意, 在忽略挠的意义下, 关系 $V=U^I$ 也表示 $V^{\perp}$ 为元素 $gx-x$ 生成的模, 其中 $x\in U$ 且 $g\in I$; 而无论如何都有

\begin{equation}\tag{\SGARefTarget{sga7.sga7i.eix.n9-2-1}{9.2.1}}
gx-x\in V^{\perp}\quad\text{当 }x\in U,\ g\in I \qquad .
\end{equation}
因而对 $g\in I$,

\begin{equation}\tag{\SGARefTarget{sga7.sga7i.eix.n9-2-2}{9.2.2}}
gx=x+u(g)\mathbin{.}x \qquad ,
\end{equation}
其中 $g\longmapsto u(g)$ 给出一个从 $U$ 到 $V^{\perp}$、在 $V=U^I$ 上为零的同态, 或等价地

\begin{equation}\tag{\SGARefTarget{sga7.sga7i.eix.n9-2-3}{9.2.3}}
u\in\operatorname{Hom}(U/V,V^{\perp})
 \simeq (U/V)^{\vee}\otimes V^{\perp} \qquad .
\end{equation}
在当前的半稳定约化情形, 即 $V^{\perp}\subset V$ 时, 公式 (\SGARefLink{sga7.sga7i.eix.n9-2-2}{9.2.2}) 表明, $g$ 在 $U$ 上的作用是平行于 $V^{\perp}=W$ 的\underline{剪切变换 (transvection)}, 而这些剪切变换的复合, 就是 (\SGARefLink{sga7.sga7i.eix.n9-2-3}{9.2.3}) 中相应同态 $u$ 的相加. 因而惯性群 $I$ 在 $U$ 上的作用, 完全由一个群同态描述

\begin{equation}\tag{\SGARefTarget{sga7.sga7i.eix.n9-2-4}{9.2.4}}
I\longrightarrow (U/V)^{\vee}\otimes V^{\perp} \qquad .
\end{equation}

分别将正交性定理用于 $A_K$ 与 $A'_K$, 得到

\begin{equation*}
V^{\perp}=W \quad , \qquad
(U/V)^{\vee}\simeq W'(-1) \qquad ,
\end{equation*}
其中 $W'$ 为 $U'=T_{\ell}(A'(\overline K))$ 的环面部分, $(-1)$ 表示 Tate 扭曲. 此外, 同态 (\SGARefLink{sga7.sga7i.eix.n9-2-4}{9.2.4}) 通过 $I(\ell)\simeq\mathbb{Z}_{\ell}(1)$ 分解, 因而 (\SGARefLink{sga7.sga7i.eix.n9-2-4}{9.2.4}) 可解释为典范同态

\begin{equation}\tag{\SGARefTarget{sga7.sga7i.eix.n9-2-5}{9.2.5}}
u:\mathbb{Z}_{\ell}(1)\longrightarrow W'\otimes W(-1)
 \qquad ,
\end{equation}
或等价地, 为一个元素

\begin{equation}\tag{\SGARefTarget{sga7.sga7i.eix.n9-2-6}{9.2.6}}
u\in W'\otimes W(-2) \qquad .
\end{equation}
回顾 $W$ 与 $W'$ 是两个扭常值 $\ell$ 进层 $T_{\ell}(A^o)^t$ 与 $T_{\ell}({A'}^o)^t$ 的 $\overline K$-值点群, 其特殊纤维分别为 $T_{\ell}(T_o)$ 与 $T_{\ell}(T'_o)$. 得到典范同构

\begin{equation}\tag{\SGARefTarget{sga7.sga7i.eix.n9-2-7}{9.2.7}}
\left\{
\begin{array}{l}
W\simeq T_{\ell}(T_o)(\overline k)
 \simeq\check M\otimes\mathbb{Z}_{\ell}(1),\\[3ex]
W'\simeq T_{\ell}(T'_o)(\overline k)
 \simeq\check M'\otimes\mathbb{Z}_{\ell}(1),
\end{array}
\right.
\end{equation}
其中 $M$ 与 $M'$ 分别表示 $T_o$ 与 $T'_o$ 在 $\overline k$ 上的特征群:

\begin{equation}\tag{\SGARefTarget{sga7.sga7i.eix.n9-2-8}{9.2.8}}
T_o=D(M) \quad , \qquad T'_o=D(M') \qquad .
\end{equation}
因此, (\SGARefLink{sga7.sga7i.eix.n9-2-6}{9.2.6}) 中描述惯性群 $I$ 在 $A_K$ 的 Tate 模 $U$ 上的作用的元素 $u$, 可解释为 $(\check M'\otimes_{\mathbb{Z}}\check M)
 \otimes_{\mathbb{Z}}\mathbb{Z}_{\ell}$ 的一个典范元素. 此外, 它在 $\pi_o=\pi/I=\operatorname{Gal}(\overline k/k)$ 下不变, 因为按定义它在结构搬移下不变:

\begin{equation}\tag{\SGARefTarget{sga7.sga7i.eix.n9-2-9}{9.2.9}}
u\in
\bigl((\check M'\otimes_{\mathbb{Z}}\check M)
 \otimes_{\mathbb{Z}}\mathbb{Z}_{\ell}\bigr)^I
\simeq
\operatorname{Hom}_{\mathbb{Z}}
 (\underline M\otimes\underline M',\mathbb{Z}_{\ell})
\simeq
\operatorname{Hom}_{\mathbb{Z}_{\ell}}
 (\underline M_{\ell}\otimes\underline M'_{\ell},
  \mathbb{Z}_{\ell}) \quad .
\end{equation}
以 $u$ 代替 $u_{\ell}$ 记之, 这就是我们要定义的元素 (\SGARefLink{sga7.sga7i.eix.n9-1-2}{9.1.2}).

为在一般情形中定义 (\SGARefLink{sga7.sga7i.eix.n9-1-2}{9.1.2}), 还需要一些预备知识, 在接下来的两节中展开.

% END INCLUDED SGA7I_English_current_buildable_source_20260823/components/158_expose_IX_section_9_2.tex

% BEGIN INCLUDED SGA7I_English_current_buildable_source_20260823/components/159_expose_IX_section_9_3.tex
% Expose IX, section 9.3.
% French transcription lines 2279--2430; authority scan physical pp. 433--439
% (scan indices 432--438; source folios 421--427). The authority pages were
% checked directly in overlapping 1100-dpi bands. Higher-detail inspection was
% reserved for the functor diagram. Evident print defects are corrected and
% disclosed in reader notes: the P/Q order in the graded pieces; F for E twice
% in the Brauer-pushout paragraph; part (d) for part (c); Ext^2(R,Q) for
% Ext^2(Q,P); omega for omega-prime on the lower functor arrow; and Y for Y'
% in the change-of-base formula.

\medskip
\noindent
\SGARefTarget{sga7.sga7i.eix.n9-3}{9.3}. \underline{混合扩张.}

设 $\underline C$ 为 阿贝尔范畴,

\begin{equation}\tag{\SGARefTarget{sga7.sga7i.eix.n9-3-1}{9.3.1}}
P,\ R,\ Q
\end{equation}
为 $\underline C$ 的三个对象. 我们要分类 $\underline C$ 中配备三步滤过的对象 $X$,

\begin{equation}\tag{*}
X^0=X\supset X^1\supset X^2\supset X^3=0
\end{equation}
并给定同构\footnote{印本依次将 $P,R,Q$ 分配给这三个分次商. 紧接着的正合列与随后所有构造都要求 $Q,R,P$; 因而此处改正了 $P$ 与 $Q$ 的明显互换.}

\begin{equation*}
X^0/X^1\xrightarrow{\ \sim\ }Q,
\qquad X^1/X^2\xrightarrow{\ \sim\ }R,
\qquad X^2/X^3=X^2\xrightarrow{\ \sim\ }P.
\end{equation*}
置

\begin{equation*}
E=X^0/X^2,
\qquad F=X^1/X^3=X^1,
\end{equation*}
可见 $X$ 给出两个通常的扩张

\begin{equation}\tag{\SGARefTarget{sga7.sga7i.eix.n9-3-2}{9.3.2}}
0\longrightarrow P\longrightarrow F\longrightarrow R\longrightarrow0,
\qquad
0\longrightarrow R\longrightarrow E\longrightarrow Q\longrightarrow0,
\end{equation}
分别为 $R$ 被 $P$ 的扩张, 以及 $Q$ 被 $R$ 的扩张. 此处的观点是, 除了对象 (\SGARefLink{sga7.sga7i.eix.n9-3-1}{9.3.1}) 外, 还从同样假定为给定的扩张结构 (\SGARefLink{sga7.sga7i.eix.n9-3-2}{9.3.2}) 出发, 考虑相应的三步滤过结构 $X$.

更精确地, 扩张 $E$ 被扩张 $F$ 的一个\underline{混合扩张 (mixed extension)}, 指 $\underline C$ 的一个对象 $X$, 配备三步滤过 (*) 以及给定的同构

\begin{equation}\tag{\SGARefTarget{sga7.sga7i.eix.n9-3-3}{9.3.3}}
E\simeq X^0/X^2,
\qquad F\simeq X^1/X^3=X^1,
\end{equation}
这些同构与滤过相容: 它们分别将 $X^1/X^2$ 映到 $\operatorname{Im}(R\longrightarrow E)$, 将 $X^2$ 映到 $\operatorname{Im}(P\longrightarrow F)$, 并要求由 (\SGARefLink{sga7.sga7i.eix.n9-3-3}{9.3.3}) 得到的同构 $R\simeq F/P$ (从而 $R\simeq X^1/X^2$), 等于由 (\SGARefLink{sga7.sga7i.eix.n9-3-2}{9.3.2}) 中第一个正合列得到的同构.

对固定的 $P,Q,R$ 与固定的扩张 $E$ 和 $F$, $E$ 被 $F$ 的混合扩张以显然方式组成一个范畴; 它事实上为群胚 (每个箭头均可逆). $E$ 被 $F$ 的混合扩张 $X$ 的自同构群典范同构于 $\operatorname{Hom}(Q,P)$, 同构由下式给出

\begin{equation}\tag{\SGARefTarget{sga7.sga7i.eix.n9-3-4}{9.3.4}}
\operatorname{Hom}(Q,P)
 \xrightarrow{\ \sim\ }
\operatorname{Aut}_{\mathrm{extpan}}(X),
\qquad
u\longmapsto\operatorname{id}_X+\bar u,
\end{equation}
其中 $\bar u$ 是 $X$ 的底层 $\underline C$ 对象的自同态, 由复合 $X\longrightarrow Q\xrightarrow{u}P\longrightarrow X$ 得到; 外侧两个箭头分别为 $X$ 上给定混合扩张结构所给出的典范满态射与典范单态射.

可以把 $E$ 被 $F$ 的混合扩张的范畴 $\operatorname{EXTPAN}(E,F)$ 的结构刻画得更精确, 即把它看作 $Q$ 被 $P$ 的扩张所成的「群范畴」(VII \SGARefLink{sga7.sga7i.evii.n2-5}{2.5}) $\operatorname{EXT}(Q,P)$ 的一个「挠子」; 所用的是与通常扩张的 Brauer 复合类似的自然函子:

\begin{equation}\tag{\SGARefTarget{sga7.sga7i.eix.n9-3-5}{9.3.5}}
\begin{aligned}
\omega:
\operatorname{EXT}(Q,P)\times\operatorname{EXTPAN}(E,F)
 &\longrightarrow \operatorname{EXTPAN}(E,F),\\
\omega(G,X)&=G\wedge X.
\end{aligned}
\end{equation}
对 $\operatorname{EXTPAN}(E,F)$ 的每个对象 $X$, 部分函子 $G\longmapsto G\wedge X$ 为从 $\operatorname{EXT}(Q,P)$ 到 $\operatorname{EXTPAN}(E,F)$ 的等价.

只简要描述 $\omega$. 将 $E$ 被 $F$ 的混合扩张 $X$ 看作 $Q$ 被 $F$ (视为 $\underline C$ 的对象) 的一个通常扩张, 仍记为 $X$. 对 $Q$ 被 $P$ 的每个扩张 $G$, 以 $\bar G$ 表示由包含态射 $P\longrightarrow F$ 诱导的 $Q$ 被 $F$ 的扩张. 置

\begin{equation}\tag{\SGARefTarget{sga7.sga7i.eix.n9-3-6}{9.3.6}}
G\wedge X=\omega(G,X)
 \overset{\mathrm{dfn}}{=}\bar G\wedge X,
\end{equation}
其中 $\wedge$ 表示 $Q$ 被 $F$ 的扩张范畴中的 Brauer 复合. (\SGARefLink{sga7.sga7i.eix.n9-3-6}{9.3.6}) 右端 $X'$ 的三步滤过显然可得, 附属于 $X'$ 的 (\SGARefLink{sga7.sga7i.eix.n9-3-3}{9.3.3}) 型的第二个同构也是如此.

为对 $X'$ 定义 (\SGARefLink{sga7.sga7i.eix.n9-3-3}{9.3.3}) 型的第一个同构, 注意函子

\begin{equation*}
Y\longmapsto\widetilde Y,
\qquad
\operatorname{EXT}(Q,F)\longrightarrow\operatorname{EXT}(Q,R),
\end{equation*}
由 (\SGARefLink{sga7.sga7i.eix.n9-3-2}{9.3.2}) 中给定的态射 $F\longrightarrow R$ 诱导, 它与 Brauer 复合相容. 因而它将 $X'$ 映为一个典范同构于 $\widetilde{\bar G}\wedge\widetilde X$ 的扩张 $\widetilde X'$. 第一个同构 (\SGARefLink{sga7.sga7i.eix.n9-3-3}{9.3.3}) 给出同构 $\widetilde X\simeq E$, 而 $\widetilde{\bar G}$, 即由 $G$ 通过复合 $P\longrightarrow F\longrightarrow R$ 诱导的扩张, 典范同构于平凡扩张, 因为此复合 $P\longrightarrow R$ 为零. 因而有典范同构 $\widetilde X'\simeq E$, 它在 (\SGARefLink{sga7.sga7i.eix.n9-3-6}{9.3.6}) 的右端规定了 $E$ 被 $F$ 的混合扩张结构.\footnote{本段中的两处, 印本均写作 $F$, 而所涉对象为 $Q$ 被 $R$ 的扩张, 因而必须与 $E$ 等同.}

将 $\omega$ 在 $\operatorname{EXT}(Q,P)\times\operatorname{EXTPAN}(E,F)$ 的箭头上的定义, (\SGARefLink{sga7.sga7i.eix.n9-3-5}{9.3.5}) 所述性质的验证 (即对每个\underline{固定}的混合扩张 $X$, 函子 $G\longmapsto G\wedge X$ 为从 $\operatorname{EXT}(Q,P)$ 到 $\operatorname{EXTPAN}(E,F)$ 的等价), 以及结合律公式的证明留给读者

\begin{equation}\tag{\SGARefTarget{sga7.sga7i.eix.n9-3-7}{9.3.7}}
G_1\wedge(G_2\wedge X)
 \simeq (G_1\wedge G_2)\wedge X,
\end{equation}
其中 $G_1\wedge G_2$ 表示 $\operatorname{EXT}(Q,P)$ 中的 Brauer 复合.

上述讨论立即给出以下结果的 (a) 与 (b).

\medskip
\noindent
\underline{命题 \SGARefTarget{sga7.sga7i.eix.n9-3-8}{9.3.8}. 在 阿贝尔范畴 $\underline C$ 中给定如 (\SGARefLink{sga7.sga7i.eix.n9-3-2}{9.3.2}) 中的两个扩张 $E$ 与 $F$, 考虑 $E$ 被 $F$ 的混合扩张范畴 $\operatorname{EXTPAN}(E,F)$. 则:}

\smallskip
\noindent
a) \underline{$\operatorname{EXTPAN}(E,F)$ 是群胚. 它的每个对象的自同构群典范同构于} $\operatorname{Ext}^0(Q,P)=\operatorname{Hom}(Q,P)$.

\smallskip
\noindent
b) \underline{$\operatorname{EXTPAN}(E,F)$ 的对象的同构类集合 $\operatorname{Extpan}(E,F)$ 或为空集, 或自然为群} $\operatorname{Ext}^1(Q,P)$ \underline{的一个挠子.}

\smallskip
\noindent
c) \underline{可以典范地定义一个类}

\begin{equation}\tag{\SGARefTarget{sga7.sga7i.eix.n9-3-9}{9.3.9}}
c(E,F)\in\operatorname{Ext}^2(Q,P),
\end{equation}
\underline{其为零的充要条件是} $\operatorname{Extpan}(E,F)\neq\varnothing$, \underline{即存在 $E$ 被 $F$ 的混合扩张.}

尚需证明 (c).\footnote{印本写作「证明 (d)」, 但命题 \SGARefLink{sga7.sga7i.eix.n9-3-8}{9.3.8} 仅有 (a)、(b)、(c) 三项.} $E$ 被 $F$ 的混合扩张也可描述为 $E$ 被 $P$ 的一个扩张 $X$, 配备一个 $R$ 被 $P$ 的扩张的同构

\begin{equation*}
X^1\simeq F,
\end{equation*}
其中 $X^1$ 表示扩张 $X$ 在 (\SGARefLink{sga7.sga7i.eix.n9-3-2}{9.3.2}) 中态射 $R\longrightarrow E$ 下的像. 因而此混合扩张存在, 当且仅当 $\operatorname{Ext}^1(E,P)$ 中存在一个元素, 其在 $\operatorname{Ext}^1(R,P)$ 中的像为 $F$ 的类 $\xi$. 考虑与 (\SGARefLink{sga7.sga7i.eix.n9-3-2}{9.3.2}) 中第二个正合列关联的 $\operatorname{Ext}^i(-,P)$ 正合列\footnote{印本的末项为 $\operatorname{Ext}^2(R,Q)$; 将 $\operatorname{Ext}(-,P)$ 用于 $0\to R\to E\to Q\to0$ 就得到所示的 $\operatorname{Ext}^2(Q,P)$, 它也是 (\SGARefLink{sga7.sga7i.eix.n9-3-9}{9.3.9}) 中指定的群.}:

\begin{equation*}
\cdots\longrightarrow
\operatorname{Ext}^1(E,P)\longrightarrow
\operatorname{Ext}^1(R,P)
\xrightarrow{\partial}
\operatorname{Ext}^2(Q,P)\longrightarrow\cdots .
\end{equation*}
只需由下式定义 $c(E,F)$

\begin{equation}\tag{\SGARefTarget{sga7.sga7i.eix.n9-3-10}{9.3.10}}
c(E,F)=\partial\xi,
\qquad
\xi=\operatorname{cl}(F)\in\operatorname{Ext}^1(R,P).
\end{equation}

现在假设给定一个正合函子

\begin{equation}\tag{\SGARefTarget{sga7.sga7i.eix.n9-3-11}{9.3.11}}
r:\underline C\longrightarrow\underline C',
\end{equation}
其中 $\underline C'$ 为另一个 阿贝尔范畴. 以 $P',Q',R',E',F'$ 表示将 $r$ 用于 $P,Q,R,E,F$ 得到的 $\underline C'$ 中的对象; 因而 $E'$ 为 $Q'$ 被 $R'$ 的扩张, $F'$ 为 $R'$ 被 $P'$ 的扩张. 精确到典范同构, 下列函子图交换:

\begin{equation}\tag{\SGARefTarget{sga7.sga7i.eix.n9-3-12}{9.3.12}}
\begin{tikzcd}[column sep=large,row sep=large]
\operatorname{EXT}(Q,P)\times\operatorname{EXTPAN}(E,F)
  \arrow[r,"\omega"]\arrow[d]
& \operatorname{EXTPAN}(E,F)\arrow[d]\\
\operatorname{EXT}(Q',P')\times\operatorname{EXTPAN}(E',F')
  \arrow[r,"\omega'"']
& \operatorname{EXTPAN}(E',F').
\end{tikzcd}
\end{equation}
竖直箭头为正合函子 $r$ 诱导的显然函子, $\omega'$ 是将 (\SGARefLink{sga7.sga7i.eix.n9-3-5}{9.3.5}) 中的 $\underline C$ 换为 $\underline C'$ 后与 $\omega$ 类似的箭头. 印本图中的下方箭头标作 $\omega$, 没有撇号; 说明句将其指明为 $\omega'$, 此处采用该记号. \SGARefLink{sga7.sga7i.eix.n9-3-8}{9.3.8} 中考虑的结构也与函子 $r$ 有类似的相容性; 其明确表述留给读者.

将使用这一交换性的下述推论. 假设 $\operatorname{EXTPAN}(E,F)$ 非空, 从中选取一个对象 $X$. 则通过 $G'\longmapsto G'\wedge X'$, 其中 $X'=r(X)$, 范畴 $\operatorname{EXTPAN}(E',F')$ 等价于 $\operatorname{EXT}(Q',P')$. 因而只要知道 $\operatorname{EXT}(Q',P')$ 中对应对象的同构类, 就精确到同构知道 $\operatorname{EXTPAN}(E',F')$ 的对象 $Y'$, 即只需知道相应元素

\begin{equation*}
c_X(Y')\in\operatorname{Ext}^1(Q',P')
\end{equation*}
即可.

若将 $X$ 换为另一个对象 $X_1=H\wedge X$, 其中 $H$ 为 $Q$ 被 $P$ 的扩张, 则立即由结合律公式得到\footnote{下式的印本右端为 $c_X(Y)$; 此构造中的固定对象为 $Y'$.}

\begin{equation*}
c_{X_1}(Y')=c_X(Y')-\operatorname{cl}(H')
\end{equation*}
即由 (\SGARefLink{sga7.sga7i.eix.n9-3-7}{9.3.7}) 得到. 因而对固定的 $Y'$, 当 $X$ 变化时, 元素 $c_X(Y')$ 恰遍历一个陪集

\begin{equation}\tag{\SGARefTarget{sga7.sga7i.eix.n9-3-13}{9.3.13}}
c(Y')\in
\operatorname{Ext}^1(Q',P')/
\operatorname{Im}\operatorname{Ext}^1(Q,P),
\end{equation}
因而它典范地关联于 $E'$ 被 $F'$ 的混合扩张 $Y'$. 由构造, 其为零的充要条件为: $Y'$ 同构于 $E$ 被 $F$ 的某个混合扩张 $X$ 在 $r$ 下的变换.

% END INCLUDED SGA7I_English_current_buildable_source_20260823/components/159_expose_IX_section_9_3.tex

% BEGIN INCLUDED SGA7I_English_current_buildable_source_20260823/components/160_expose_IX_section_9_4.tex
% Expose IX, section 9.4.
% French transcription lines 2435--2576; authority scan physical pp. 439--444
% (scan indices 438--443; source folios 427--432). The authority pages were
% checked directly in overlapping 1100-dpi bands. The scan is clear throughout;
% no higher-resolution escalation was required. Evident print/transcription
% defects are corrected and disclosed at the point of use: omitted checks and
% delimiters, Lambda_s misread as a wedge, Q/P reversals in Hom and Ext, and the
% nonexistent reference 9.4.6.2 for 9.4.7.2.

\medskip
\noindent
\SGARefTarget{sga7.sga7i.eix.n9-4}{9.4}. 设 $S$ 为一般点为 $\eta$ 的 Hensel 性状, $n$ 为整数 $>0$, 置 $\Lambda=\mathbb Z/n\mathbb Z$, 令 $\underline C$ 为 $S$ 上的 $\Lambda$-模范畴 (对于 $S$ 上的 fppf 景; 参见 VIII \SGARefLink{sga7.sga7i.eviii.n0-2}{0.2}), $\underline C'$ 为 $\eta$ 上的 $\Lambda$-模范畴. 设 $P,R,Q$ 为 $\underline C$ 的对象, $E,F$ 分别为 $Q$ 被 $R$ 的扩张以及 $R$ 被 $P$ 的扩张 (\SGARefLink{sga7.sga7i.eix.n9-3-2}{9.3.2}). 假设 $Q$ 是有限型局部自由 $\Lambda_S$-模, $P$ 由 $S$ 上的有限乘法型群表示 (SGA 3 IX), 因而它形如

\begin{equation*}
P=D(Q^*),
\end{equation*}
其中 $D=\underline{\operatorname{Hom}}(-,\mathbb G_{mS})$ 表示 Cartier 对偶. 以 $P_\eta,R_\eta,Q_\eta,E_\eta,F_\eta$ 表示 $P,R,Q,E,F$ 到 $\eta$ 的限制, 并假设给定 $E_\eta$ 被 $F_\eta$ 的一个混合扩张 $X_\eta$. 在这些条件下, 将典范地定义一个配对

\begin{equation}\tag{\SGARefTarget{sga7.sga7i.eix.n9-4-1}{9.4.1}}
c(X_\eta):Q\otimes_\Lambda Q^*\longrightarrow\Lambda_S
\qquad(\Lambda=\mathbb Z/n\mathbb Z).
\end{equation}

为此首先\underline{假设} $S$ \underline{严格局部}. 则 $Q$ 与 $Q^*$ 为常值群, 设分别为 $Q=Q_{oS}$ 与 $Q^*=Q^*_{oS}$, 而给出 (\SGARefLink{sga7.sga7i.eix.n9-4-1}{9.4.1}) 等价于给出一个元素

\begin{equation}\tag{\SGARefTarget{sga7.sga7i.eix.n9-4-2}{9.4.2}}
c(X_\eta)\in\check Q_o\otimes\check Q_o^*,
\end{equation}
其中倒帽记号表示通常的对偶 $\operatorname{Hom}(-,\Lambda)$. 对限制函子 $r:\underline C\longrightarrow\underline C'$ 应用 \SGARefLink{sga7.sga7i.eix.n9-3}{9.3} 的讨论. 首先注意, 存在 $E$ 被 $F$ 的混合扩张. 事实上, 其存在性的障碍属于 $\operatorname{Ext}^2(Q,P)$ (\SGARefLink{sga7.sga7i.eix.n9-3-8}{9.3.8}(c)), 因而只需证明

\begin{equation*}
\operatorname{Ext}^2(Q,P)=0.
\end{equation*}

这是下列引理的特例.

\medskip
\noindent
\underline{引理 \SGARefTarget{sga7.sga7i.eix.n9-4-3}{9.4.3}. 在上述假设下, 若 $S$ 严格局部, 则有:}

\smallskip
\noindent
a) \underline{典范同构}

\begin{equation*}
\operatorname{Ext}^i_\Lambda(Q,P)\simeq
H^i(S,\check Q\otimes P),
\qquad
\operatorname{Ext}^i_\Lambda(Q_\eta,P_\eta)\simeq
H^i(\eta,\check Q\otimes P);
\end{equation*}

\noindent
b) \underline{典范同构}

\begin{equation}\tag{b)}
\check Q\otimes P\simeq D(Q\otimes Q^*);
\end{equation}

\noindent
c) \underline{对 $S$ 上每个有限常值群 $H=H_{oS}$,}

\begin{equation}\tag{c)}
H^i(S,D(H))\simeq H^i(\eta,D(H)_\eta)=0
\qquad\underline{\text{当 }}i\geq2;
\end{equation}

\noindent
d) \underline{典范同构}

\begin{equation}\tag{d)}
\begin{aligned}
H^1(S,D(H))&\simeq\mathbb G_m(S)\otimes\check H_o
              =V^*\otimes\check H_o,\\
H^1(\eta,D(H)_\eta)&\simeq\mathbb G_m(\eta)\otimes\check H_o
              =K^*\otimes\check H_o,
\end{aligned}
\end{equation}
\underline{其中 (d) 中的倒帽记号表示 $\operatorname{Hom}(-,\mathbb
Q/\mathbb Z)$.}

(a) 和 (b) 中的公式立即来自 $Q$ 为有限型局部自由 $\Lambda_S$-模这一事实. 为证明 (c), 显然可归结到 $H_o=\mathbb Z/m\mathbb Z$ 的情形, 即 $D(H)=\mu_m$. 与 Kummer 正合列关联的上同调列

\begin{equation*}
0\longrightarrow\mu_m\longrightarrow\mathbb G_m
 \xrightarrow{\ m\cdot\operatorname{id}\ }\mathbb G_m
 \longrightarrow0
\end{equation*}
将问题归结为证明

\begin{equation*}
H^i(S,\mathbb G_m)=H^i(\eta,\mathbb G_m)=0
\qquad\text{当 }i\geq1.
\end{equation*}
这些群是在平展景上计算的群 $H^i_{\mathrm{\acute et}}$ [6, 定理 \SGAExternalRef{sga7.sga7i.eix.n11-7}{11.7}{文献 6 的编号，非本报告内部编号}], 因而当 $S$ 时消没, 因为 $S$ 严格局部. 它们在 $\eta$ 时的消没也众所周知 [6, 定理 \SGAExternalRef{sga7.sga7i.eix.n1-1}{1.1}{文献 6 的编号，非本报告内部编号}].

尚需描述 (d) 中的同构. 在任意底 $U$ 上, 存在典范同态

\begin{equation}\tag{\SGARefTarget{sga7.sga7i.eix.n9-4-3-1}{9.4.3.1}}
\mathbb G_m(U)\otimes\check H_o
\longrightarrow H^1(U,D(H_{oU})).
\end{equation}
若 $H_o$ 被整数 $m>0$ 零化, 则此同态可用 Kummer 列关联的边缘映射 $\mathbb G_m(U)\xrightarrow{\partial}H^1(U,\mu_m)$ 描述: 它是下列复合

\begin{equation*}
\begin{aligned}
\mathbb G_m(U)\otimes\check H_o
&\longrightarrow H^1(U,\mu_m)\otimes\check H_o\\
&\longrightarrow H^1(U,\mu_m\otimes\check H_o)
 \longrightarrow H^1(U,D(H_{oU})).
\end{aligned}
\end{equation*}
其中第二个箭头由显然的同态 $\check H_o\longrightarrow
\operatorname{Hom}(\mu_m,\mu_m\otimes\check H_o)$ 得到, 第三个是 (b) 中以 $m$ 代替 $n$ 的特例. 留给读者验证, (\SGARefLink{sga7.sga7i.eix.n9-4-3-1}{9.4.3.1}) 与满足 $mH_o=0$ 的整数 $m$ 的选取无关. 还断言, (\SGARefLink{sga7.sga7i.eix.n9-4-3-1}{9.4.3.1}) 总是单射, 且当 $\operatorname{Pic}(U)=0$ 时其余核为零. (事实上, 余核总同构于 $\operatorname{Hom}(H_o,\operatorname{Pic}(U))$, 但这里无关紧要.) 事实上, 归结到 $H_o=\mathbb Z/n\mathbb Z$ 的情形, 该断言立即由上述 Kummer 列的上同调列得到. 这就证明了 (d), 并立即给出:

\medskip
\noindent
\underline{推论 \SGARefTarget{sga7.sga7i.eix.n9-4-4}{9.4.4}. 沿用 \SGARefLink{sga7.sga7i.eix.n9-4-3}{9.4.3} 的记号, 存在典范同构}

\noindent
e)
\begin{equation*}
H^1(\eta,D(H)_\eta)/\operatorname{Im}H^1(S,D(H))
\simeq\check H_o,
\end{equation*}
\underline{从而结合 (a) 与 (b), 得到典范同构}

\begin{equation}\tag{\SGARefTarget{sga7.sga7i.eix.n9-4-5}{9.4.5}}
\operatorname{Ext}^1(Q_\eta,P_\eta)/
\operatorname{Im}\operatorname{Ext}^1(Q,P)
\simeq\check Q_o\otimes\check Q_o^*.
\end{equation}

回到 $S$ 严格局部时 (\SGARefLink{sga7.sga7i.eix.n9-4-2}{9.4.2}) 的定义, 关系 $\operatorname{Ext}^2(Q,P)=0$ 使我们能够取 $Y'$ 为 $E_\eta$ 被 $F_\eta$ 的给定混合扩张 $X_\eta$, 定义典范元素 (\SGARefLink{sga7.sga7i.eix.n9-3-13}{9.3.13}). 公式 (\SGARefLink{sga7.sga7i.eix.n9-4-5}{9.4.5}) 将此元素 $c(X_\eta)$ 等同于 $\check Q_o\otimes\check Q_o^*$ 的元素, 从而完成 (\SGARefLink{sga7.sga7i.eix.n9-4-2}{9.4.2}) 的定义. \SGARefLink{sga7.sga7i.eix.n9-3}{9.3} 的一般讨论还表明, 此元素就是将 $X_\eta$ 延拓为 $E$ 被 $F$ 的混合扩张 $X$ 的障碍.

\medskip
\noindent
\SGARefTarget{sga7.sga7i.eix.n9-4-6}{9.4.6}. 现在当 $S$ 为 Hensel 性状时, 从 $S$ 的严格 Hensel 化作 Galois 下降, 立即将 (\SGARefLink{sga7.sga7i.eix.n9-4-1}{9.4.1}) 的定义归结到严格局部情形.

\medskip
\noindent
\underline{命题 \SGARefTarget{sga7.sga7i.eix.n9-4-7}{9.4.7}. 设 $P,Q,R,E,F$ 如上, 但性状 $S$ 不必为 Hensel 性状. 考虑函子}

\begin{equation}\tag{\SGARefTarget{sga7.sga7i.eix.n9-4-7-1}{9.4.7.1}}
X\longmapsto X_\eta
\end{equation}
\underline{它从 $E$ 被 $F$ 的混合扩张范畴到 $E_\eta$ 被 $F_\eta$ 的混合扩张范畴. 对 $E_\eta$ 被 $F_\eta$ 的每个混合扩张 $X_\eta$, 考虑配对}

\begin{equation}\tag{\SGARefTarget{sga7.sga7i.eix.n9-4-7-2}{9.4.7.2}}
c_o(X_\eta):Q_o\otimes Q_o^*\longrightarrow\Lambda_s
\end{equation}
\underline{它由 $S$ 的 Hensel 化 $S^h$ 上定义的 (\SGARefLink{sga7.sga7i.eix.n9-4-1}{9.4.1}) 型配对诱导. 则函子 (\SGARefLink{sga7.sga7i.eix.n9-4-7-1}{9.4.7.1}) 全忠实, 而 $E_\eta$ 被 $F_\eta$ 的混合扩张 $X_\eta$ 属于其本质像, 当且仅当配对 (\SGARefLink{sga7.sga7i.eix.n9-4-7-2}{9.4.7.2}) 为零.}

(\SGARefLink{sga7.sga7i.eix.n9-4-7-2}{9.4.7.2}) 在印本中的目标为 $\Lambda_s$; 法文转录中的 $\wedge_s$ 为文字识别错误.

\medskip
\noindent
\SGARefTarget{sga7.sga7i.eix.n9-4-7-3}{9.4.7.3}. 若只假设 $S$ 为维数为 $1$ 的连通正则诺特概形, 此陈述仍有意义且成立. 此时应对 $S$ 的每个闭点 $s$ 考虑一个配对 $Q_s\otimes Q_s^*\longrightarrow\Lambda_s$, 而不再只有单个配对 (\SGARefLink{sga7.sga7i.eix.n9-4-7-2}{9.4.7.2}). 在这一更一般形式下证明 \SGARefLink{sga7.sga7i.eix.n9-4-7}{9.4.7} 最为方便.

\smallskip
\noindent
a) 对每个混合扩张 $X$, 映射

\begin{equation*}
\operatorname{End}(X)\longrightarrow\operatorname{End}(X_\eta)
\end{equation*}
是双射. 两边分别为 $\underline{\operatorname{Hom}}(Q,P)=D(Q\otimes Q^*)$ 在 $S$ 与 $\eta$ 上的截面; 因而该断言由 $S$ 的正规性以及 $D(Q\otimes Q^*)$ 在 $S$ 上的有限性得到.\footnote{印本证明此处写作 $\underline{\operatorname{Hom}}(P,Q)$. 命题 9.3.8(a)、公式 (\SGARefLink{sga7.sga7i.eix.n9-3-4}{9.3.4}) 以及所示的与 $D(Q\otimes Q^*)$ 的等同都要求 $\underline{\operatorname{Hom}}(Q,P)$.} 立即推出, 若 $X$ 与 $Y$ 为已知同构的两个混合扩张, 则

\begin{equation*}
\operatorname{Hom}(X,Y)\longrightarrow
\operatorname{Hom}(X_\eta,Y_\eta)
\end{equation*}
是双射.

\smallskip
\noindent
b) 若 $X$ 与 $Y$ 为两个混合扩张, 则每个同构 $X_\eta\xrightarrow{\sim}Y_\eta$ 来自一个同构 $X\longrightarrow Y$. 由 (a), 这一问题对平展拓扑是局部的, 因而可归结到 $S$ 严格局部的情形. 此时 \SGARefLink{sga7.sga7i.eix.n9-4-3}{9.4.3} 表明映射

\begin{equation*}
\operatorname{Ext}^1_\Lambda(Q,P)
\longrightarrow\operatorname{Ext}^1_\Lambda(Q_\eta,P_\eta)
\end{equation*}
等同于

\begin{equation*}
V^*\otimes\check H_o\longrightarrow K^*\otimes\check H_o,
\qquad H_o=Q_o\otimes Q_o^*.
\end{equation*}
后者为单射, 因为正合列 $0\longrightarrow V^*\longrightarrow K^*\longrightarrow\mathbb Z
\longrightarrow0$ 分裂. 由 \SGARefLink{sga7.sga7i.eix.n9-3}{9.3} 的讨论, 这蕴涵若 $X_\eta$ 与 $Y_\eta$ 同构, 则 $X$ 与 $Y$ 也同构; 再由 (a) 即得结论.\footnote{印本在本段中各 Ext 群的参数为 $(P,Q)$. 命题 \SGARefLink{sga7.sga7i.eix.n9-3-8}{9.3.8}(b) 中 $\operatorname{Ext}^1(Q,P)$ 在混合扩张上的作用, 以及引理 \SGARefLink{sga7.sga7i.eix.n9-4-3}{9.4.3} 本身, 都要求 $(Q,P)$.}

因而 (a) 与 (b) 表明 $X\longmapsto X_\eta$ 是\underline{全}忠实的.

\smallskip
\noindent
c) 设 $X_\eta$ 为 $\eta$ 上的混合扩张. 证明它可延拓到 $S$, 当且仅当配对 (\SGARefLink{sga7.sga7i.eix.n9-4-7-2}{9.4.7.2}) 消没. 这里同样, 已证的唯一性表明该问题对平展拓扑是局部的, 再次归结到严格局部情形. 如 \SGARefLink{sga7.sga7i.eix.n9-4-6}{9.4.6} 之前所指出, 此时结论按构造即成立.\footnote{印本证明引用了不存在的配对「(9.4.6.2)」; 应为配对 (\SGARefLink{sga7.sga7i.eix.n9-4-7-2}{9.4.7.2}).}

这就完成了 \SGARefLink{sga7.sga7i.eix.n9-4-7}{9.4.7} 的证明.

% END INCLUDED SGA7I_English_current_buildable_source_20260823/components/160_expose_IX_section_9_4.tex

% BEGIN INCLUDED SGA7I_English_current_buildable_source_20260823/components/161_expose_IX_section_9_5.tex
% Expose IX, section 9.5.
% French transcription lines 2578--2653; authority scan physical pp. 444--448
% (scan indices 443--447; source folios 432--436). The authority pages were
% checked directly in four overlapping 1100-dpi bands. The scan is clear
% throughout. Four internally forced print slips are corrected and disclosed:
% 9.4.6 for 9.4.7, X_eta for its extension X over S, eta for S as the base of
% the etale group Q(n), and 9.5.6 for the homomorphism (9.5.5).

\medskip
\noindent
\SGARefTarget{sga7.sga7i.eix.n9-5}{9.5}. 固定素数 $\ell$, $S$ 上的 pro-$\ell$ Barsotti--Tate 群

\begin{equation}\tag{\SGARefTarget{sga7.sga7i.eix.n9-5-1}{9.5.1}}
P=(P(n))_{n\geq0},
\qquad R=(R(n))_{n\geq0},
\qquad Q=(Q(n))_{n\geq0},
\end{equation}
以及 pro-$\ell$ Barsotti--Tate 群的扩张

\begin{equation}\tag{\SGARefTarget{sga7.sga7i.eix.n9-5-2}{9.5.2}}
0\longrightarrow P\longrightarrow F\longrightarrow R\longrightarrow0,
\qquad
0\longrightarrow R\longrightarrow E\longrightarrow Q\longrightarrow0.
\end{equation}
因而对每个整数 $n\geq0$, 它们定义 $\Lambda_n$-模的扩张, 其中 $\Lambda_n=\mathbb Z/\ell^{\,n+1}\mathbb Z$:

\begin{equation}\tag{\SGARefTarget{sga7.sga7i.eix.n9-5-3}{9.5.3}}
\begin{gathered}
0\longrightarrow P(n)\longrightarrow F(n)\longrightarrow R(n)
 \longrightarrow0,\\
0\longrightarrow R(n)\longrightarrow E(n)\longrightarrow Q(n)
 \longrightarrow0.
\end{gathered}
\end{equation}
假设 $Q$ 为平展型, $P$ 为乘法型. 最后, 假设给定 $\eta$ 上的 pro-$\ell$ Barsotti--Tate 群 $X_\eta$, 它是 $E_\eta$ 被 $F_\eta$ 的\underline{混合扩张}. 这意味着它作为 pro-$\ell$ Barsotti--Tate 群配备三步滤过, 各部分扩张分别配备到 $E_\eta$ 与 $F_\eta$ 的同构. 因而对每个 $n$, $X_\eta(n)$ 具有 $E(n)_\eta$ 被 $F(n)_\eta$ 的混合扩张结构. 于是可给它对应一个 (\SGARefLink{sga7.sga7i.eix.n9-4-1}{9.4.1}) 型配对:

\begin{equation}\tag{\SGARefTarget{sga7.sga7i.eix.n9-5-4}{9.5.4}}
c(X(n)_\eta):Q(n)\otimes Q^*(n)
\longrightarrow\Lambda_n=\mathbb Z/\ell^{\,n+1}\mathbb Z,
\end{equation}
其中各 $Q^*(n)$ 为由下式定义的 $S$ 上的有限平展群

\begin{equation*}
P(n)\simeq D(Q^*(n)),
\qquad\text{即}\qquad Q^*(n)=D(P(n)),
\end{equation*}
而 $D$ 仍表示 Cartier 对偶. 各 $Q^*(n)$ 以显然方式组成 pro-$\ell$ 群, 记为 $Q^*$. 留给读者验证, 当 $n$ 变化时, 各同态 (\SGARefLink{sga7.sga7i.eix.n9-5-4}{9.5.4}) 定义投射系统的同态, 即 $S$ 上扭常值 $\ell$ 进层的同态:

\begin{equation}\tag{\SGARefTarget{sga7.sga7i.eix.n9-5-5}{9.5.5}}
c(X_\eta):Q\otimes Q^*\longrightarrow\mathbb Z_{\ell S}.
\end{equation}

顺便提及下列结果, 下文不会使用它.

\medskip
\noindent
\underline{命题 \SGARefTarget{sga7.sga7i.eix.n9-5-6}{9.5.6}. 设 $P,Q,R,E,F$ 如上, 但性状 $S$ 不必为 Hensel 性状. 考虑函子}

\begin{equation}\tag{\SGARefTarget{sga7.sga7i.eix.n9-5-6-1}{9.5.6.1}}
X\longmapsto X_\eta
\end{equation}
\underline{它从 $E$ 被 $F$ 的 Barsotti--Tate 混合扩张范畴到 $E_\eta$ 被 $F_\eta$ 的混合扩张范畴. 对 $E_\eta$ 被 $F_\eta$ 的混合扩张 Barsotti--Tate 群 $X_\eta$, 考虑相应配对}

\begin{equation}\tag{\SGARefTarget{sga7.sga7i.eix.n9-5-6-2}{9.5.6.2}}
Q_o\otimes Q_o^*\longrightarrow\mathbb Z_{\ell s}
\end{equation}
\underline{它由 $S$ 的 Hensel 化上定义的 (\SGARefLink{sga7.sga7i.eix.n9-5-5}{9.5.5}) 型配对诱导. 则函子 (\SGARefLink{sga7.sga7i.eix.n9-5-6-1}{9.5.6.1}) 全忠实, 其本质像恰由以下 Barsotti--Tate 群 $X_\eta$ 组成: 它们为 $E_\eta$ 被 $F_\eta$ 的混合扩张, 且对应配对 (\SGARefLink{sga7.sga7i.eix.n9-5-6-2}{9.5.6.2}) 为零. 此外, 当 $\operatorname{char}K=0$ 时 (从而 Tate 定理 [33, 定理 4] 适用), $X_\eta$ 属于本质像, 当且仅当它可以延拓为 $S$ 上的 Barsotti--Tate 群.}

\medskip
\noindent
\SGARefTarget{sga7.sga7i.eix.n9-5-6-3}{9.5.6.3}. 留给读者将此陈述推广到仅假设 $S$ 为维数为 $1$ 的连通正则诺特概形的情形.

\SGARefLink{sga7.sga7i.eix.n9-5-6}{9.5.6} 的第一个断言, 立即由 \SGARefLink{sga7.sga7i.eix.n9-4-7}{9.4.7} 应用于所考虑 Barsotti--Tate 群的各分量得到.\footnote{印本证明引用 \SGARefLink{sga7.sga7i.eix.n9-4-6}{9.4.6}, 它仅进行 Hensel 下降构造; 逐分量应用的全忠实性与本质像陈述为命题 \SGARefLink{sga7.sga7i.eix.n9-4-7}{9.4.7}.} 尚需证明: 若 $E_\eta$ 被 $F_\eta$ 的混合扩张 Barsotti--Tate 群 $X_\eta$ 可延拓为 $S$ 上的 Barsotti--Tate 群 $X$, 则它也可作为混合扩张延拓.\footnote{在此, 印本同样将 $S$ 上的延拓记为 $X_\eta$. 此处保留 $X_\eta$ 表示一般纤维, 以 $X$ 表示其在 $S$ 上的延拓.} 由 Tate 定理 [33], 同态 $X_\eta\longrightarrow Q_\eta$ 唯一地延拓为同态

\begin{equation*}
X\longrightarrow Q.
\end{equation*}
对每个整数 $n$, 相应同态 $X(n)\longrightarrow Q(n)$ 在一般纤维上为满射, 因而是满射. 由于 $Q(n)$ 在 $S$ 上平展, 映射 $X(n)\longrightarrow Q(n)$ 平坦, 从而为满态射. 因此 $X\longrightarrow Q$ 是\underline{逐分量的满态射}.\footnote{印本此句称 $Q(n)$ 在 $\eta$ 上平展. 群 $Q(n)$ 是在 $S$ 上定义的, 并在该处被假设为平展型; 这一平坦性论证所需的底为 $S$.} 令 $\overline F$ 为 $X\longrightarrow Q$ 的核. 则 $\overline F$ 为 Barsotti--Tate 群, $X$ 为 $Q$ 被 $\overline F$ 的扩张. 同构 $F_\eta\xrightarrow{\sim}\overline F_\eta$ 由 Tate 定理延拓为同构 $F\simeq\overline F$. 因而 $X$ 为 $Q$ 被 $F$ 的扩张. 现在考虑商 $\overline E=X/P$. 再次应用 Tate 定理, 得到同构 $E\simeq\overline E$. 因而 $X$ 为 $E$ 被 $F$ 的混合扩张, 并延拓混合扩张 $X_\eta$, 如所要求.

\medskip
\noindent
\underline{注 \SGARefTarget{sga7.sga7i.eix.n9-5-6-4}{9.5.6.4}.}

\smallskip
\noindent
a) 当然, \SGARefLink{sga7.sga7i.eix.n9-5-6}{9.5.6} 的最后一个断言应当不受 $K$ 的限制而成立, 因为所引用的 Tate 定理也应如此.

\smallskip
\noindent
b) 可以对好约化判据 \SGARefLink{sga7.sga7i.eix.n5-10}{5.10} 给出另一证明. 如同 \SGARefLink{sga7.sga7i.eix.n5-10}{5.10}, 归结到 $A_K$ 有半稳定约化的情形, 再将 \SGARefLink{sga7.sga7i.eix.n9-5-6}{9.5.6} 用于视为混合扩张的 $T_\ell(A_K)$ (参见 \SGARefLink{sga7.sga7i.eix.n9-6}{9.6}). 这蕴涵 $u_\ell=0$; 再利用下面的 \SGARefLink{sga7.sga7i.eix.n11-6-2}{11.6.2}(a), 得到 $T_\ell(A_K)^t=0$.

\medskip
\noindent
\SGARefTarget{sga7.sga7i.eix.n9-5-8}{9.5.8}. 当 $\ell\neq p$ 时, 分别在 $S$ 与 $K$ 上考虑的 Barsotti--Tate 群, 可以直接解释为 $S$ 上的扭常值 $\ell$ 进层, 或等价地, 分别解释为 $\pi_o=\operatorname{Gal}(\overline k/k)$ 与 $\pi=\operatorname{Gal}(\overline K/K)$ 在有限型自由 $\mathbb Z_\ell$-模上的表示. 于是可见, 同态 (\SGARefLink{sga7.sga7i.eix.n9-5-5}{9.5.5}) 就是由 \SGARefLink{sga7.sga7i.eix.n9-2}{9.2} 的方法定义的、表达惯性子群 $I\subset\pi$ 在 $X_\eta(\overline K)$ 上的作用的典范同态.\footnote{印本此处引用「同态 (\SGARefLink{sga7.sga7i.eix.n9-5-6}{9.5.6})」; 本节构造的同态是 (\SGARefLink{sga7.sga7i.eix.n9-5-5}{9.5.5}), 而 \SGARefLink{sga7.sga7i.eix.n9-5-6}{9.5.6} 是命题.} 为验证这一点, 立即归结到 \SGARefLink{sga7.sga7i.eix.n9-4}{9.4} 的情形, 在该处取 $n$ 为 $\ell$ 的幂, $R$ 为局部常值 $\Lambda$-模 (如同 $Q$ 和 $P$); 于是 $E,F$ 也如此, 且 $X_\eta$. 两个定义一致的验证留给读者.

% END INCLUDED SGA7I_English_current_buildable_source_20260823/components/161_expose_IX_section_9_5.tex

% BEGIN INCLUDED SGA7I_English_current_buildable_source_20260823/components/162_expose_IX_sections_9_6_and_9_7.tex
% Expose IX, sections 9.6--9.7.
% French transcription lines 2655--2711; authority scan physical pp. 448--450
% (scan indices 447--449; source folios 436--438). The authority pages were
% checked directly in four overlapping 1100-dpi bands. The scan is clear.
% Four evident print slips affecting the English are corrected and disclosed:
% F_eta for E_eta in the quotient extension, (9.5.6) for (9.5.5), S for eta
% as the base of X_eta, and a missing closing parenthesis. A missing sentence
% stop in 9.7 is restored.

\medskip
\noindent
\SGARefTarget{sga7.sga7i.eix.n9-6}{9.6}. 现在可以不对 $\ell$ 作任何限制, 定义预告的典范配对 (\SGARefLink{sga7.sga7i.eix.n9-1-2}{9.1.2}), 并使其在 $\ell\neq p$ 时与 \SGARefLink{sga7.sga7i.eix.n9-2}{9.2} 的定义一致. 为此将 \SGARefLink{sga7.sga7i.eix.n9-5}{9.5} 的讨论用于如下情形:

\begin{equation}\tag{\SGARefTarget{sga7.sga7i.eix.n9-6-1}{9.6.1}}
\left\{
\begin{aligned}
P&=T_\ell(A^o)^t=\check{\underline M}{}'_\ell(1),
& Q&=\underline M_\ell,
& R&=T_\ell(A^o)^f/T_\ell(A^o)^t,\\[1.25ex]
F&=T_\ell(A^o)^f,
& E&=D(F')=D\bigl(T_\ell({A'}^o)^f\bigr).
\end{aligned}
\right.
\end{equation}
最后一个公式中, $D$ 表示 pro-$\ell$ Barsotti--Tate 群意义下的 Cartier 对偶; 即 $D(Y)$ 为各 $D(Y(n))$ 配以通常过渡态射所组成的\underline{投射}系统. $F$ 作为 $R$ 被 $P$ 的扩张的结构是显然的. 类似地,

\begin{equation*}
F'=T_\ell({A'}^o)^f
\end{equation*}
具有 $R'$ 被 $P'=T_\ell({A'}^o)^t=\check{\underline M}_\ell(1)$ 的扩张结构. 应用 Cartier 对偶, 得到 $E=D(F')$ 具有 $D(P')=\underline M_\ell$ 被 $D(R')$ 的扩张结构. 而 $D(R')$ 由 \SGARefLink{sga7.sga7i.eix.n5-5}{5.5} (在 \SGARefLink{sga7.sga7i.eix.n7-4-6}{7.4.6} 中证明) 典范同构于 $R$; 这就规定了 $E$ 作为 $Q$ 被 $R$ 的扩张的结构.

对于 $E_\eta$ 被 $F_\eta$ 的混合扩张 $X_\eta$, 取

\begin{equation}\tag{\SGARefTarget{sga7.sga7i.eix.n9-6-2}{9.6.2}}
X_\eta=T_\ell(A_K)=X^0_\eta
 \supset X^1_\eta=T_\ell(A_K)^f
 \supset X^2_\eta=T_\ell(A_K)^t
 \supset0.
\end{equation}
以下 $R_\eta$ 被 $P_\eta$ 的扩张的同构

\begin{equation*}
X^1_\eta\simeq F_\eta
\end{equation*}
是显然的. 以下 $Q_\eta$ 被 $R_\eta$ 的扩张的同构

\begin{equation*}
X_\eta/X^2_\eta\simeq E_\eta=D(F'_\eta)
\end{equation*}
由正交性定理 \SGARefLink{sga7.sga7i.eix.n5-2}{5.2} 显然可得.\footnote{印本公式写作 $F_\eta=D(F'_\eta)$. 由于 $E=D(F')$, 且所示商是 $Q_\eta$ 被 $R_\eta$ 的扩张, 所需的项为 $E_\eta$.}

因而我们具备了定义配对 (\SGARefLink{sga7.sga7i.eix.n9-5-5}{9.5.5}) 所需的资料, 这里因 $Q^*=\underline M'_\ell$, 配对具有 (\SGARefLink{sga7.sga7i.eix.n9-1-2}{9.1.2}) 的形式.\footnote{印本称其为配对「(\SGARefLink{sga7.sga7i.eix.n9-5-6}{9.5.6})」; 此配对为同态 (\SGARefLink{sga7.sga7i.eix.n9-5-5}{9.5.5}), 而 \SGARefLink{sga7.sga7i.eix.n9-5-6}{9.5.6} 是随后的命题.} 它与 \SGARefLink{sga7.sga7i.eix.n9-2}{9.2} 的定义的相容性来自 \SGARefLink{sga7.sga7i.eix.n9-5-8}{9.5.8}. \SGARefLink{sga7.sga7i.eix.n5-8}{5.8} 的讨论还表明, 在 Tate 定理 [33] 有效的条件下 (因而至少当 $\operatorname{char}K=0$ 时), 此配对可以只由 $\eta$ 上单个 pro-$\ell$ Barsotti--Tate 群 $X_\eta=T_\ell(A_K)$ 描述: 其三步滤过以及延拓 $X^1_\eta$ 和 $X_\eta/X^2_\eta$ 的扩张 $F$ 与 $E$, 均可仅由 $X_\eta$ 典范地定义.\footnote{印本此句称 $X_\eta$「在 $S$ 上」; 其记号与定义表明它是 $\eta$ 上的一般纤维群. 此处印本还缺少右括号.}

\medskip
\noindent
\underline{注 \SGARefTarget{sga7.sga7i.eix.n9-7}{9.7}.} 不再要求 $A_K$ 在 Hensel 性状 $S$ 上必有半稳定约化. 令 $K'$ 为 $K$ 的一个有限 Galois 扩张, 使 $A_{K'}$ 在 $S$ 于 $K'$ 中的正规化 $S'$ 上有半稳定约化 (\SGARefLink{sga7.sga7i.eix.n3-6}{3.6}). 沿用 (\SGARefLink{sga7.sga7i.eix.n8-1-8}{8.1.8}) 的记号, 对 $A_{K'}$ 的单值配对作 Galois 下降, 便定义一个配对

\begin{equation}\tag{*}
v_\ell^{K'}:\underline M_K\otimes\underline M'_K
 \longrightarrow\mathbb Z_{\ell K}.
\end{equation}
先验而言, 此配对依赖于 $K'$ 的选取, 但由下面的 (10.4.3), 配对

\begin{equation}\tag{\SGARefTarget{sga7.sga7i.eix.n9-7-1}{9.7.1}}
u_\ell:\underline M_K\otimes\underline M'_K
 \longrightarrow\mathbb Q_{\ell K}
\end{equation}
由下式给出

\begin{equation*}
u_\ell=\frac{1}{e(S'/S)}\,v_\ell^{K'}
\end{equation*}
与该选取无关, 其中

\begin{equation*}
e(S'/S)=\operatorname{length}(V'/\mathfrak mV').
\end{equation*}
仍可将 (\SGARefLink{sga7.sga7i.eix.n9-7-1}{9.7.1}) 称为 $K$ 上阿贝尔概形 $A_K$ 对当前所考虑素数 $\ell$ 的\underline{单值配对}: 当 $A_K$ 本身有半稳定约化时, 显然它与单值配对 (\SGARefLink{sga7.sga7i.eix.n9-1-2}{9.1.2}) 到 $\eta$ 的限制一致. 由下面的 \SGARefLink{sga7.sga7i.eix.n10-4}{10.4} 可知, 事实上此配对总取值于 $\mathbb Q_K$, 且与 $\ell$ 无关.

可以问它是否取值于 $\mathbb Z_K$, 这等价于说, 对每个素数 $\ell$, (\SGARefLink{sga7.sga7i.eix.n9-7-1}{9.7.1}) 都取值于 $\mathbb Z_{\ell K}$. 至少当 $A_K$ 为椭圆曲线时确实如此. 此时唯一需验证的是 $\underline M_K\simeq\underline M'_K$ 的秩为 $1$ 的情形; 于是 $\underline M_K\otimes\underline M'_K$ 典范同构于 $\mathbb Z_K$, 单值配对 (\SGARefLink{sga7.sga7i.eix.n9-7-1}{9.7.1}) 可直接解释为一个通常的有理数. 由 Néron 的结果 [19], 此数为 $v(j)$, 其中 $j$ 为 $A_K$ 的绝对不变量, $v$ 为 $K$ 的正规化赋值.


% END INCLUDED SGA7I_English_current_buildable_source_20260823/components/162_expose_IX_sections_9_6_and_9_7.tex

% BEGIN INCLUDED SGA7I_English_current_buildable_source_20260823/components/163_expose_IX_sections_10_0_through_10_2.tex
% Expose IX, section 10 opening and sections 10.1--10.2.
% French transcription lines 2717--2786; authority scan physical pp. 451--453
% (scan indices 450--452; source folios 439--441). Physical p. 454 was
% checked for the exact 10.3 boundary. All pages were compared in four
% overlapping 1100-dpi bands; the print is clear at that resolution.

\medskip
\noindent
10. \underline{单值配对的性质}. \underline{整性与正性定理}

本节首先给出单值配对的一些实质上平凡的形式性质, 将验证的 (稍显繁琐的) 细节留给读者. 还将陈述关于这些配对的一个更深的定理 (\SGARefLink{sga7.sga7i.eix.n10-4}{10.4}), 并说明如何将其证明归结到 Jacobian 簇的一个特例; 此特例将在第 12 节中 (用超越方法) 处理, 该节在逻辑上事实上独立于第 11 节.

\medskip
\noindent
\SGARefTarget{sga7.sga7i.eix.n10-1}{10.1}. \underline{结构搬移}. 由单值配对 (\SGARefLink{sga7.sga7i.eix.n9-1-2}{9.1.2}) 的定义显然可知, 此配对与情形 $(S,A_K,A'_K)$ 中的结构搬移相容.

\medskip
\noindent
\SGARefTarget{sga7.sga7i.eix.n10-2}{10.2}. \underline{对称性}. 从 $A_K$ 的对偶概形 $A'_K$ 出发, 对它应用第 9 节的构造, 必须引入 $A'_K$ 的对偶概形 $(A'_K)'$, 它由双对偶定理 ([5] 或 VIII 3.2) 典范同构于 $A_K$. 以 $u_\ell^{A_K}$ 表示关于 $A_K$ 的单值配对 (\SGARefLink{sga7.sga7i.eix.n9-1-2}{9.1.2}). 同样得到配对

\begin{equation}\tag{\SGARefTarget{sga7.sga7i.eix.n10-2-1}{10.2.1}}
u_\ell^{A'_K}:\underline M'_\ell\otimes\underline M_\ell
 \longrightarrow\mathbb Z_\ell.
\end{equation}
断言后者由 $u_\ell^{A_K}$ 通过对称同构

\begin{equation*}
s:\underline M'_\ell\otimes\underline M_\ell
 \xrightarrow{\ \sim\ }
 \underline M_\ell\otimes\underline M'_\ell;
\end{equation*}
得到, 即

\begin{equation}\tag{\SGARefTarget{sga7.sga7i.eix.n10-2-2}{10.2.2}}
u_\ell^{A'_K}=u_\ell^{A_K}\circ s.
\end{equation}

考虑配对 $\varphi_\xi$ 已知的反交换性质 ([5] 或 VIII 2.2.11), 可以利用 $A_K$ 的固定极化 $\xi$, 将对称公式 (\SGARefLink{sga7.sga7i.eix.n10-2-2}{10.2.2}) 改写如下. 它定义同态

\begin{equation*}
\widetilde\xi:A_K\longrightarrow A'_K
\end{equation*}
从而定义相应同态 (其中 $t$ 表示「环面」)

\begin{equation}\tag{\SGARefTarget{sga7.sga7i.eix.n10-2-3}{10.2.3}}
\xi_t:\underline M\longrightarrow\underline M'.
\end{equation}
因此可以形成配对

\begin{equation}\tag{\SGARefTarget{sga7.sga7i.eix.n10-2-4}{10.2.4}}
\begin{aligned}
\varphi^t_{\xi,\ell}:\underline M_\ell\otimes\underline M_\ell
 &\longrightarrow\mathbb Z_\ell,\\
\varphi^t_{\xi,\ell}(x,y)&=u_\ell^{A_K}\bigl(x,\xi_t(y)\bigr).
\end{aligned}
\end{equation}
上述型 $\varphi^t_{\xi,\ell}$ 是\underline{对称的}:

\begin{equation}\tag{\SGARefTarget{sga7.sga7i.eix.n10-2-5}{10.2.5}}
\varphi^t_{\xi,\ell}(x,y)=\varphi^t_{\xi,\ell}(y,x).
\end{equation}
事实上, 此断言对 $A_K$ 的 Néron--Severi 群的每个元素 $\xi$ 都成立, 并不一定要求为极化. 但只有在 $\xi$ 非退化时, 即它定义一个\underline{同源} $A_K\longrightarrow A'_K$, 从而也定义同源 (\SGARefLink{sga7.sga7i.eix.n10-2-3}{10.2.3}) 时, 对称关系 (\SGARefLink{sga7.sga7i.eix.n10-2-5}{10.2.5}) 才等价于 (\SGARefLink{sga7.sga7i.eix.n10-2-2}{10.2.2}). 当 $\xi$ 为极化时, 下面的 \SGARefLink{sga7.sga7i.eix.n10-4}{10.4}(b) 还将大幅加强 (\SGARefLink{sga7.sga7i.eix.n10-2-5}{10.2.5}).

最后提及对称关系的第三种解释. 考虑 $K$ 上具有半稳定约化的一对阿贝尔概形 $(A_K,\overline A_K)$ 上的一个双扩张 (即除子对应) $\xi$. 它定义同态

\begin{equation*}
A_K\longrightarrow\overline A'_K.
\end{equation*}
转到 Néron 模型后, 它在特殊纤维的极大环面上定义同态, 从而在相反方向上定义其特征群之间的同态:

\begin{equation}\tag{\SGARefTarget{sga7.sga7i.eix.n10-2-6}{10.2.6}}
\xi_t:\overline{\underline M}\longrightarrow\underline M',
\end{equation}
其中 $\overline{\underline M}$ 由 $\overline A_K$ 定义的方式, 与 $\underline M$ 由 $A_K$ 定义的方式相同. 因而可以形成配对

\begin{equation}\tag{\SGARefTarget{sga7.sga7i.eix.n10-2-7}{10.2.7}}
\begin{aligned}
\varphi^t_\xi:\underline M_\ell\otimes
 \overline{\underline M}_\ell&\longrightarrow\mathbb Z_\ell,\\
\varphi^t_\xi(x,y)&=u_\ell^{A_K}\bigl(x,\xi_t(y)\bigr).
\end{aligned}
\end{equation}
\footnote{印本右端的底层配对仅记为 $\varphi^t$, 没有给出定义. 它是相同构造 (\SGARefLink{sga7.sga7i.eix.n10-2-4}{10.2.4}) 中使用的单值配对 $u_\ell^{A_K}$.} 顺便注意, 对变动的具有半稳定约化的阿贝尔概形 $A_K$, 扭常值群 $\underline M$ 是 $A_K$ 的协变函子. 因而, 与 $(K,S)$ 上具有半稳定约化的阿贝尔概形的双扩张关联的配对 (\SGARefLink{sga7.sga7i.eix.n10-2-7}{10.2.7}) 的构成, 在显然意义下与双扩张的逆像相容.

当 $\xi$ 为 $(A_K,A'_K)$ 的典范 (「Weil」) 双扩张时, 配对 (\SGARefLink{sga7.sga7i.eix.n10-2-7}{10.2.7}) 当然恰为单值配对 (\SGARefLink{sga7.sga7i.eix.n9-1-2}{9.1.2}).\footnote{印本此处引用 (\SGARefLink{sga7.sga7i.eix.n9-2-1}{9.2.1}), 它是惯性作用的正交公式 $gx-x\in V^\perp$, 并不是单值配对. 此配对的编号为 (\SGARefLink{sga7.sga7i.eix.n9-1-2}{9.1.2}).} 现在考虑 $\xi$ 的转置 ${}^s\xi$, 它是 $(\overline A_K,A_K)$ 的双扩张, 因而给出配对

\begin{equation*}
\varphi^t_{{}^s\xi}:\overline{\underline M}_\ell\otimes
 \underline M_\ell\longrightarrow\mathbb Z_\ell.
\end{equation*}
对称关系为

\begin{equation}\tag{\SGARefTarget{sga7.sga7i.eix.n10-2-8}{10.2.8}}
\varphi^t_{{}^s\xi}=\varphi^t_\xi\circ s,
\qquad
s:\overline{\underline M}_\ell\otimes\underline M_\ell
 \xrightarrow{\ \sim\ }
 \underline M_\ell\otimes\overline{\underline M}_\ell,
\end{equation}
其中 $s$ 仍为对称同构. 可将此关系与 VIII \SGARefLink{sga7.sga7i.eviii.n2-2-11}{2.2.11} 的反对称关系比较.


% END INCLUDED SGA7I_English_current_buildable_source_20260823/components/163_expose_IX_sections_10_0_through_10_2.tex

% BEGIN INCLUDED SGA7I_English_current_buildable_source_20260823/components/164_expose_IX_sections_10_3_and_10_4.tex
% Expose IX, sections 10.3--10.4.
% French transcription lines 2792--2883; authority scan physical pp. 454--456
% (scan indices 453--455; source folios 442--444). Physical p. 457 was
% checked for the exact 10.5 boundary. All pages were compared in four
% overlapping 1100-dpi bands; the text, formulas, and simple square are
% clear at that resolution.

\medskip
\noindent
\SGARefTarget{sga7.sga7i.eix.n10-3}{10.3}. \underline{性状变换}. 设

\begin{equation}\tag{\SGARefTarget{sga7.sga7i.eix.n10-3-1}{10.3.1}}
g:S'\longrightarrow S
\end{equation}
为 Hensel 性状的控制态射, 对应于离散赋值环的单的局部同态

\begin{equation*}
V\longrightarrow V'.
\end{equation*}
如常置

\begin{equation}\tag{\SGARefTarget{sga7.sga7i.eix.n10-3-2}{10.3.2}}
e(S'/S)=\operatorname{length}_V\bigl(V'/\underline mV'\bigr),
\end{equation}
其中 $\underline m$ 为 $V$ 的极大理想. 已知 (\SGARefLink{sga7.sga7i.eix.n10-3-1}{10.3.1}) 在惯性群的「温和」部分上诱导的同态,

\begin{equation}\tag{\SGARefTarget{sga7.sga7i.eix.n10-3-3}{10.3.3}}
g_*:I'_t\simeq\prod_{\ell\ne p}\mathbb Z_\ell(1)(k')
 \longrightarrow
 I_t=\prod_{\ell\ne p}\mathbb Z_\ell(1)(k),
\end{equation}
由下式给出

\begin{equation}\tag{\SGARefTarget{sga7.sga7i.eix.n10-3-4}{10.3.4}}
g_*(x)=e(S'/S)x,
\end{equation}
这里利用显然的同构 $\mathbb Z_\ell(1)(k)\xrightarrow{\ \sim\ }\mathbb Z_\ell(1)(k')$ 等同了 (\SGARefLink{sga7.sga7i.eix.n10-3-3}{10.3.3}) 的两边.\footnote{印本在「报告 I」之后的内部引用留白.}

另一方面, 若 $A_K$ 为 $K$ 上在 $S$ 上有半稳定约化的阿贝尔概形, 则扭格 $\underline M$ 与 $\underline M'$ 的构成与上述任意基变换 $S'\longrightarrow S$ 交换, 这立即来自定义 (特别是 \SGARefLink{sga7.sga7i.eix.n3-4}{3.4}). 因而可以将 $A_{K'}$ 定义的配对 $u_\ell^{A_{K'}}$ 看作下述配对, 其中 $K'$ 为 $V'$ 的分式域,

\begin{equation*}
u_\ell^{A_{K'}}:
 \underline M_{\ell S'}\otimes\underline M'_{\ell S'}
 \longrightarrow\mathbb Z_{\ell S'},
\end{equation*}
并与由 $u_\ell^{A_K}$ 经基变换得到的配对 $u_{\ell S'}^{A_K}$ 比较. 当 $\ell\ne p$ 时, \SGARefLink{sga7.sga7i.eix.n9-2}{9.2} 中配对 $u_\ell$ 的构造结合公式 (\SGARefLink{sga7.sga7i.eix.n10-3-4}{10.3.4}), 立即表明\footnote{印本引用 (\SGARefLink{sga7.sga7i.eix.n10-3-3}{10.3.3}), 即所示的映射 $g_*$; 此处使用的乘法公式为 (\SGARefLink{sga7.sga7i.eix.n10-3-4}{10.3.4}).}

\begin{equation}\tag{\SGARefTarget{sga7.sga7i.eix.n10-3-5}{10.3.5}}
u_\ell^{A_{K'}}=e(S'/S)u_\ell^{A_K}.
\end{equation}
事实上此公式对 $\ell=p$ 仍然成立. 为此回到 \SGARefLink{sga7.sga7i.eix.n9-3}{9.3}--\SGARefLink{sga7.sga7i.eix.n9-6}{9.6} 的一般构造, 归结到严格局部性状变换 $S'\longrightarrow S$ 下同构 \SGARefLink{sga7.sga7i.eix.n9-4-4}{9.4.4}(e) 的类似相容性断言. 注意此同构由正规化赋值 $K^*/V^*\simeq\mathbb Z$ 诱导, 因而给出交换图\footnote{印本称 $K^*/V^*\simeq\mathbb Z$ 由「范数」得到. 所示映射为正规化的离散赋值.}

\begin{center}
\begin{tikzcd}
K^*/V^* \arrow[r, "\sim"] \arrow[d]
  & \mathbb Z \arrow[d, "e(S'/S)"] \\
K^{\prime *}/V^{\prime *} \arrow[r, "\sim"]
  & \mathbb Z .
\end{tikzcd}
\end{center}

下面给出关于单值配对的一个更深的结果. 将在第 12 节用超越方法证明它, 所依据的结果将在讨论班后续报告中证明. (事实上, 在讨论班其余部分中不会使用 \SGARefLink{sga7.sga7i.eix.n10-4}{10.4}, 如同不会使用本报告的其他结果一样.)

\medskip
\noindent
\underline{定理 \SGARefTarget{sga7.sga7i.eix.n10-4}{10.4}.} \underline{设 $S$ 为 Hensel 性状, $A_K$ 为 $K$ 上在 $S$ 上有半稳定约化的阿贝尔概形. 考虑如 \SGARefLink{sga7.sga7i.eix.n9-1}{9.1} (参见 \SGARefLink{sga7.sga7i.eix.n5-5-2}{5.5.2}) 中关联的格 $\underline M$ 与 $\underline M'$. 则下列断言成立:}

\smallskip
\noindent
a) \underline{存在唯一的配对}

\begin{equation}\tag{\SGARefTarget{sga7.sga7i.eix.n10-4-1}{10.4.1}}
u:\underline M\otimes\underline M'\longrightarrow\mathbb Z_S
\end{equation}
\underline{使得对每个素数 $\ell$, $u\otimes_{\mathbb Z}\mathbb Z_\ell$ 就是 (\SGARefLink{sga7.sga7i.eix.n9-1-2}{9.1.2}) 中的单值配对 $u_\ell$. 此配对是分离的.}

\smallskip
\noindent
b) \underline{设 $\xi$ 为 $A_K$ 的一个极化, 从而定义同态 (\SGARefLink{sga7.sga7i.eix.n10-2-3}{10.2.3})}

\begin{equation*}
\xi_t:\underline M\longrightarrow\underline M',
\end{equation*}
\underline{并考虑型}

\begin{equation}\tag{\SGARefTarget{sga7.sga7i.eix.n10-4-2}{10.4.2}}
\begin{aligned}
\varphi_\xi^t:\underline M\otimes\underline M
  &\longrightarrow\mathbb Z_S,\\
\varphi_\xi^t(x,y)&=u\bigl(x,\xi_t(y)\bigr).
\end{aligned}
\end{equation}
\underline{此型对称、正且非退化.}

当然, 若以 $S$ 的某个给定几何点处 $\underline M$ 的几何纤维 $M_o$ 来表述, 则上述断言说相应的型

\begin{equation*}
M_o\otimes M_o\longrightarrow\mathbb Z
\end{equation*}
对称、正且非退化. 配对 (\SGARefLink{sga7.sga7i.eix.n10-4-1}{10.4.1}) 仍称为\underline{单值配对}.

% END INCLUDED SGA7I_English_current_buildable_source_20260823/components/164_expose_IX_sections_10_3_and_10_4.tex

% BEGIN INCLUDED SGA7I_English_current_buildable_source_20260823/components/165_expose_IX_section_10_5.tex
% Expose IX, section 10.5.
% French transcription lines 2889--2955; authority scan physical pp. 457--460
% (scan indices 456--459; source folios 445--448). Physical p. 460 was
% checked through the exact section 11 boundary. All pages were compared in
% four overlapping 1100-dpi bands; the prose and formulas are clear at that
% resolution, so no higher-resolution escalation was necessary.

\medskip
\noindent
\SGARefTarget{sga7.sga7i.eix.n10-5}{10.5}. \underline{将 \SGARefLink{sga7.sga7i.eix.n10-4}{10.4} 归结到一个特例}

\smallskip
\noindent
\SGARefTarget{sga7.sga7i.eix.n10-5-1}{10.5.1}. 满足所述条件的 $u$ 的唯一性显然可见 (甚至只需对一个 $\ell$ 满足条件). 另一方面, 存在性仅意味着每个 $u_\ell$ 事实上将 $\underline M\otimes\underline M'$ 映到 $\mathbb Z_S$, 且由此得到的同态

\begin{equation*}
\underline M\otimes\underline M'\longrightarrow\mathbb Z_S
\end{equation*}
与 $\ell$ 无关. 事实上甚至只需证明对各 $u_\ell\otimes_{\mathbb Z_\ell}\mathbb Q_\ell$ 以及一个同态

\begin{equation*}
u_{\mathbb Q}:\underline M_{\mathbb Q}\otimes
 \underline M'_{\mathbb Q}\longrightarrow\mathbb Q_S.
\end{equation*}

成立的类似断言. 假设经过如 \SGARefLink{sga7.sga7i.eix.n9-3}{9.3} 中的基变换 $S'\longrightarrow S$ 后已证明存在性. 由 (\SGARefLink{sga7.sga7i.eix.n10-3-5}{10.3.5}) 立即得到存在同态

\begin{equation*}
u:\underline M\otimes\underline M'
 \longrightarrow\frac{1}{e(S'/S)}\mathbb Z_S\subset\mathbb Q_S
\end{equation*}
与各 $u_\ell\otimes_{\mathbb Z_\ell}\mathbb Q_\ell$ 相容; 这就证明了 a). 对于 b), 显然, 若它对 $A_{K'}$ 成立, 则也对 $A_K$ 成立.

因此, 在证明 \SGARefLink{sga7.sga7i.eix.n10-4}{10.4} 时可作任意控制性状变换 $S'\longrightarrow S$. 特别地, 这允许假设 $S$ 完备, 且剩余域代数闭. 此时 $\underline M$ 和 $\underline M'$ 为常值群概形, 取值分别为有限生成自由 $\mathbb Z$-模 $M$ 和 $M'$, 可以更方便地用后者重述 a) 与 b).

\smallskip
\noindent
\SGARefTarget{sga7.sga7i.eix.n10-5-2}{10.5.2}. 现在假设 $A_K$ 同构于 $K$ 上有半稳定约化的某个阿贝尔概形 $\overline A_K$ 的直因子, 且对后者已证明 \SGARefLink{sga7.sga7i.eix.n10-4}{10.4}. 考虑格 $\underline M,\underline M'$ 及单值配对对于具有半稳定约化的阿贝尔概形的直和的显然函子性, 立即推出 \SGARefLink{sga7.sga7i.eix.n10-4}{10.4} 对 $A_K$ 也成立.\footnote{印本称 $\overline A_K$「在 $k$ 上」; 记号以及直因子论证都要求为「在 $K$ 上」.}

同样, 由 \SGARefLink{sga7.sga7i.eix.n10-5-1}{10.5.1} 立即得到, \SGARefLink{sga7.sga7i.eix.n10-4}{10.4} 的有效性仅依赖于 $A_K$ 的同源类. 结合上述注记与「Poincar\'e 可约性定理」[15], 可知若 $A_K$ 为一个满足 \SGARefLink{sga7.sga7i.eix.n10-4}{10.4} 的阿贝尔概形 $\overline A_K$ 的商, 则 \SGARefLink{sga7.sga7i.eix.n10-4}{10.4} 对 $A_K$ 也成立.\footnote{印本在引文 [15] 后的具体位置留白.}

\smallskip
\noindent
\SGARefTarget{sga7.sga7i.eix.n10-5-3}{10.5.3}. 现在使用众所周知的事实 [31, 定理~11 (第~20~页)]: 必要时对 $K$ 作有限扩张后 (由 \SGARefLink{sga7.sga7i.eix.n10-5-1}{10.5.1} 允许), 每个阿贝尔概形 $A_K$ 都是一个 Jacobian 簇的商,

\begin{equation*}
\underline{\operatorname{Pic}}^0_{X_K/K}
\end{equation*}
它是一条 $K$ 上紧合、光滑、几何连通曲线 $X_K$ 的 Jacobian 簇. 因而由 \SGARefLink{sga7.sga7i.eix.n10-5-2}{10.5.2}, 归结到对此种 Jacobian 簇证明 \SGARefLink{sga7.sga7i.eix.n10-4}{10.4}.\footnote{印本称曲线 $X_K$「在 $X$ 上」; 所示的相对 Picard 概形以及此论证都要求为「在 $K$ 上」.}

\smallskip
\noindent
\SGARefTarget{sga7.sga7i.eix.n10-5-4}{10.5.4}. 由 Artin--Mumford 定理 [3] (已在报告 V 中引用并使用), 必要时再作 $K$ 的有限扩张 $K'$ 后, 可设 $X_K$ 延拓为 $S$ 上的射影平坦概形, 它正则, 且其几何特殊纤维 $X_{\bar s}$ 除普通二次奇点外没有其他奇点 (这里即两个分支的正规交叉).\footnote{印本在「报告 V」之后的内部引用留白.} 在此情形, 将在下面的 \SGARefLink{sga7.sga7i.eix.n12-5}{12.5} 中完全用此情形的几何资料描述单值配对 (\SGARefLink{sga7.sga7i.eix.n10-4-1}{10.4.1}), 从而证明 \SGARefLink{sga7.sga7i.eix.n10-4}{10.4} a); 同时还将对 Jacobian 簇的典范极化验证 \SGARefLink{sga7.sga7i.eix.n10-4}{10.4} b).

\smallskip
\noindent
\SGARefTarget{sga7.sga7i.eix.n10-5-5}{10.5.5}. 尚需说明: 若对给定的具有半稳定约化的 $A_K$, 断言 \SGARefLink{sga7.sga7i.eix.n10-4}{10.4} a) 成立, 且断言 \SGARefLink{sga7.sga7i.eix.n10-4}{10.4} b) 对 $A_K$ 的某个给定极化 $\xi_o$ 成立, 则同一结论对 $A_K$ 的每个极化 $\xi$ 都成立. 事实上, 设

\begin{equation*}
\widetilde\xi_o,\widetilde\xi:A_K\longrightarrow A'_K
\end{equation*}
为 $\xi_o$ 与 $\xi$ 所定义的同态; 它们为同源. 则

\begin{equation*}
\widetilde\xi=\widetilde\xi_o\alpha
 ={}^t(\alpha^*)\widetilde\xi_o,
 \qquad
\alpha\in\operatorname{End}(A_K)\otimes_{\mathbb Z}\mathbb Q,
\end{equation*}
其中 $\alpha\longmapsto\alpha^*$ 表示由极化 $\xi_o$ 在代数 $\operatorname{End}(A_K)\otimes_{\mathbb Z}\mathbb Q$ 上定义的典范对合, ${}^t(\alpha^*)$ 中的上标 $t$ 表示转置 (将 $A_K$ 的一个自同态对应为 $A'_K$ 的一个自同态). 另一方面, 已知 [15, 第~V~章, 定理~3]

\begin{equation*}
\alpha=\sum_i\beta_i^*\beta_i,
 \qquad
\beta_i\in\operatorname{End}(A_K)\otimes_{\mathbb Z}\mathbb Q,
\end{equation*}
从而

\begin{equation*}
\widetilde\xi=\sum_i{}^t\beta_i\,\widetilde\xi_o\,\beta_i.
\end{equation*}
等价地, 若用与所考虑极化关联的双扩张 (即除子对应) 来表示 (VIII \SGARefLink{sga7.sga7i.eviii.n4-14}{4.14}),

\begin{equation}\tag{\SGARefTarget{sga7.sga7i.eix.n10-5-5-1}{10.5.5.1}}
\delta(\xi)=\sum_i\delta(\xi_o)\circ(\beta_i,\beta_i).
\end{equation}
由 (\SGARefLink{sga7.sga7i.eix.n10-2-7}{10.2.7}) 之后的讨论, 此公式蕴涵 (沿用 (\SGARefLink{sga7.sga7i.eix.n10-4-2}{10.4.2}) 的记号)

\begin{equation}\tag{\SGARefTarget{sga7.sga7i.eix.n10-5-5-2}{10.5.5.2}}
\varphi^t_\xi
 =\sum_i\varphi^t_{\xi_o}(\beta_i,\beta_i)
 \in\operatorname{Hom}
 \bigl(M_{\mathbb Q}\otimes_{\mathbb Q}M_{\mathbb Q},\mathbb Q\bigr).
\end{equation}
这确实表明, 若 $\varphi^t_{\xi_o}$ 对称且正, 则 $\varphi^t_\xi$ 也如此.

最后, $\varphi^t_{\xi_o}$ 的非退化性显然等价于配对 (\SGARefLink{sga7.sga7i.eix.n10-4-1}{10.4.1}) 是分离的这一断言 (因为 $\widetilde\xi_o$ 为同源, 所以 $\xi_{ot}:M\otimes_{\mathbb Z}\mathbb Q
\longrightarrow M'\otimes_{\mathbb Z}\mathbb Q$ 是同构), 因而同样与所选极化无关. 还可注意到, 为验证 (\SGARefLink{sga7.sga7i.eix.n10-4-1}{10.4.1}) 为分离的, 只需对某个素数 $\ell$ 验证 (\SGARefLink{sga7.sga7i.eix.n9-1-2}{9.1.2}) 的单值配对 $u_\ell$ 的相同断言. 取 $\ell\ne p$, 可用 \SGARefLink{sga7.sga7i.eix.n9-2}{9.2} 中对 $u_\ell$ 的描述; 于是 $u_\ell$ 在左边为分离的 (因秩的原因从而在两边均为分离的) 这一事实立即可得.

% END INCLUDED SGA7I_English_current_buildable_source_20260823/components/165_expose_IX_section_10_5.tex

% BEGIN INCLUDED SGA7I_English_current_buildable_source_20260823/components/166_expose_IX_sections_11_0_through_11_2.tex
% Expose IX, sections 11.0--11.2.
% French transcription lines 2956--3027; authority scan physical pp. 460--463
% (scan indices 459--462; source folios 448--451). Physical p. 463 was
% checked through the exact 11.3 boundary. All pages were compared in four
% overlapping 1100-dpi bands; the text and formulas are clear at that
% resolution, so no higher-resolution escalation was necessary.

\medskip
\noindent
11. \underline{N\'eron 模型的连通分支: 对偶关系与渐近行为}

\smallskip
\noindent
\SGARefTarget{sga7.sga7i.eix.n11-0}{11.0}. 本节的目的是研究第 1 节中已考虑过的有限平展 $k$-群

\begin{equation}\tag{\SGARefTarget{sga7.sga7i.eix.n11-0-1}{11.0.1}}
\Phi_o=A_o/A_o^{,o}
\end{equation}
对每个素数 $\ell$, 以

\begin{equation*}
\Phi_o(\ell)
\end{equation*}
表示其 $\ell$-准素分量. 因而由公式可知, 知道 $\Phi_o$ 等价于知道所有 $\Phi_o(\ell)$,

\begin{equation}\tag{\SGARefTarget{sga7.sga7i.eix.n11-0-2}{11.0.2}}
\Phi_o=\coprod_\ell\Phi_o(\ell).
\end{equation}

注意, 群 $\Phi_o$ 与 $\Phi_o(\ell)$ 分别由通常的有限群 $\Phi_o(\bar k)$ 与 $\Phi_o(\ell)(\bar k)$, 连同下述群在其上的自然作用所确定,

\begin{equation*}
\pi_o=\operatorname{Gal}(\bar k/k)
\end{equation*}
另一方面, 显然 $\Phi_o(\bar k)$ 也是 $A_{\widetilde K}$ 的 N'eron 模型的连通分支群, 其中 $\widetilde K$ 为 $S$ 关于所选的 $k$ 的可分闭包 $\bar k$ 的严格 Hensel 化的分式域. 因而研究 $\Phi_o$ 时, 可归结到 $S$ 严格局部的情形; 群

\begin{equation*}
\pi_o\simeq\operatorname{Gal}(\widetilde K/K)
\end{equation*}
在 $\Phi_o(\bar k)$ 上的作用随即通过结构搬移自动恢复. 还可进一步归结到 $S$ 完备的情形, 所用的是基变换 $\widehat S\longrightarrow S$ 以及 N\'eron 模型的构成与之交换这一事实 [21, 定理~\SGAExternalRef{sga7.sga7i.eix.n3-4}{3.4}{文献 21 的编号，非本报告内部编号} b)].

\smallskip
\noindent
\SGARefTarget{sga7.sga7i.eix.n11-1}{11.1}. 因而假设 $S$ 严格局部, 于是 $\Phi_o$ 可等同于通常的有限群 $\Phi_o(k)$. 得到

\begin{equation}\tag{\SGARefTarget{sga7.sga7i.eix.n11-1-1}{11.1.1}}
\Phi_o(k)\simeq A_o(k)/A_o^o(k)
 \simeq A(S)/A^o(S)
 \simeq A(K)/A(K)^o,
\end{equation}
第一个等式来自 $A_o$ 在 $k$ 上光滑, 第二个等式来自 Hensel 引理, 因为 $A$ 在 $S$ 上光滑, 且 $S$ 为 Hensel 性状; 第三个等式来自关系 $A(S)\xrightarrow{\sim}A(K)$ 及约定

\begin{equation}\tag{\SGARefTarget{sga7.sga7i.eix.n11-1-2}{11.1.2}}
A(K)^o\mathrel{\mathop:}=A^o(S).
\end{equation}
因而 $A(K)^o$ 自然成为 $A(K)$ 中有限指数的子群 (并且不能与 $A^o(K)=A(K)$ 混淆). 同构 (\SGARefLink{sga7.sga7i.eix.n11-1-1}{11.1.1}) 显然关于 $A$ 为函子性的, 并与情形 $(S,A_K)$ 中的结构搬移相容, 特别地与任意 Galois 群的作用相容.

现在令 $\ell$ 为不同于 $p$ 的素数. 则 $A^o$ 是 $\ell$-可除的; 更精确地, $\ell\operatorname{id}_{A^o}$ 为平展满射态射. 因而它诱导 $A^o(S)=A(K)^o$ 的一个满射自同态, 所以后者是一个\underline{$\ell$-可除群}. 立即得到, 对每个素于 $p=\operatorname{char}(k)$ 的整数 $n\geq 0$,

\begin{equation}\tag{\SGARefTarget{sga7.sga7i.eix.n11-1-3}{11.1.3}}
{}_n\Phi_o(k)\simeq{}_nA(K)/{}_nA(K)^o,
 \qquad (n,p)=1.
\end{equation}
由 ${}_nA_K$ 的不变部分的定义 (\SGARefLink{sga7.sga7i.eix.n2-2-3}{2.2.3}), 有

\begin{equation}\tag{\SGARefTarget{sga7.sga7i.eix.n11-1-4}{11.1.4}}
{}_nA(K)^o=({}_nA_K)^f(\bar K)
 \simeq T_n(A_K)^f(\bar K)
 \otimes_{\mathbb Z_n}\mathbb Z/n\mathbb Z.
\end{equation}
取 $n$ 为 $\ell$ 的幂, 并如 (\SGARefLink{sga7.sga7i.eix.n2-5-3}{2.5.3}) 中置

\begin{equation}\tag{\SGARefTarget{sga7.sga7i.eix.n11-1-5}{11.1.5}}
U=T_\ell\bigl(A(\bar K)\bigr),
 \qquad
T_\ell\bigl(A(\bar K)\bigr)^f=U^I,
\end{equation}
其中 $U^I$ 表示 $U$ 中在惯性群 $I$ 下不变的子模, 可将 (\SGARefLink{sga7.sga7i.eix.n11-1-3}{11.1.3}) 写成

\begin{equation}\tag{\SGARefTarget{sga7.sga7i.eix.n11-1-6}{11.1.6}}
{}_n\Phi_o(k)\simeq (U_n)^I/(U^I)_n,
 \qquad n=\ell^{r+1},
\end{equation}
其中下标 $n$ 表示模 $n$ 约化, 即在 $\mathbb Z_\ell$ 上与 $\mathbb Z_\ell/n\mathbb Z_\ell$ 作张量积. (\SGARefLink{sga7.sga7i.eix.n11-1-6}{11.1.6}) 两边各群之间的过渡态射是显然的, 对 $\ell$ 的各幂取归纳极限, 得到下述结果.

\medskip
\noindent
\underline{命题 \SGARefTarget{sga7.sga7i.eix.n11-2}{11.2}.} \underline{假设 $S$ 严格局部, $\ell$ 为不同于 $p$ 的素数. 则存在典范同构}

\begin{equation}\tag{\SGARefTarget{sga7.sga7i.eix.n11-2-1}{11.2.1}}
\Phi_o(\ell)\simeq
 \frac{(U\otimes D_\ell)^I}{U^I\otimes D_\ell},
\end{equation}
\underline{其中 $U$ 在 (\SGARefLink{sga7.sga7i.eix.n11-1-5}{11.1.5}) 中定义, 且 $D_\ell=\mathbb Q_\ell/\mathbb Z_\ell$.}

因而群 $\Phi_o(\ell)$ 可以仅用 Tate 模 $U=T_\ell(A(\bar K))$ 来描述, 这里将它看作惯性群 $I$ 作用下的 Galois 模.

% END INCLUDED SGA7I_English_current_buildable_source_20260823/components/166_expose_IX_sections_11_0_through_11_2.tex

% BEGIN INCLUDED SGA7I_English_current_buildable_source_20260823/components/167_expose_IX_section_11_3.tex
% Expose IX, section 11.3.
% French transcription lines 3029--3147; authority scan physical pp. 463--466
% (scan indices 462--465; source folios 451--454). Physical p. 466 was
% checked through the exact 11.4 boundary. All pages were compared in four
% overlapping 1100-dpi bands. The text, formulas, arrow directions, labels,
% and two native diagrams are clear at that resolution.

\medskip
\noindent
\SGARefTarget{sga7.sga7i.eix.n11-3}{11.3}. 将用显式表达式 (\SGARefLink{sga7.sga7i.eix.n11-2-1}{11.2.1}) 概述典范配对的定义

\begin{equation}\tag{\SGARefTarget{sga7.sga7i.eix.n11-3-1}{11.3.1}}
\Phi_o(\ell)\otimes\Phi'_o(\ell)\longrightarrow\mathbb Q/\mathbb Z
\end{equation}
它是完美对偶; 这里如同第 1 节, $\Phi'_o$ 表示与 $A_K$ 的\underline{对偶} $A'_K$ 的 N\'eron 模型关联的连通分支群. 几乎可以肯定 (尽管作者没有进行验证), 此配对正是 (\SGARefLink{sga7.sga7i.eix.n1-2-1}{1.2.1}) 中定义的配对, 至多差一个符号; 请读者作为练习, 建立所需的相容性 (并确定正确的符号!).

置

\begin{equation*}
U'=T_\ell\bigl(A'(\bar K)\bigr),
 \qquad
T=T_\ell(1)(\bar K)=T_\ell(\bar K^*).
\end{equation*}
存在与 $I$ 的作用相容的有限生成自由 $\mathbb Z_\ell$-模的\underline{完美配对},

\begin{equation}\tag{\SGARefTarget{sga7.sga7i.eix.n11-3-2}{11.3.2}}
\varphi:U\otimes U'\longrightarrow T,
 \qquad (T\simeq\mathbb Z_\ell\text{ 非典范地}),
\end{equation}
并由 $I$ 的已知结构, 存在正合列\footnote{印本在 $I$ 第一次出现后的内部引用留白.}

\begin{equation}\tag{\SGARefTarget{sga7.sga7i.eix.n11-3-3}{11.3.3}}
1\longrightarrow R\longrightarrow I\longrightarrow T\longrightarrow1,
\end{equation}
其中 $R$ 为\underline{素于 $\ell$} 的投射有限群. 现在从资料 (\SGARefLink{sga7.sga7i.eix.n11-3-2}{11.3.2}) 与 (\SGARefLink{sga7.sga7i.eix.n11-3-3}{11.3.3}) 出发, 引入群

\begin{equation}\tag{\SGARefTarget{sga7.sga7i.eix.n11-3-4}{11.3.4}}
\begin{aligned}
\Phi&=\frac{(U\otimes D_\ell)^I}{U^I\otimes D_\ell},\qquad
\Phi'=\frac{(U'\otimes D_\ell)^I}{{U'}^I\otimes D_\ell},\qquad
D_\ell=\mathbb Q_\ell/\mathbb Z_\ell.
\end{aligned}
\end{equation}
将定义一个完美对偶

\begin{equation}\tag{\SGARefTarget{sga7.sga7i.eix.n11-3-5}{11.3.5}}
\Phi\otimes\Phi'\longrightarrow\mathbb Q/\mathbb Z,
\end{equation}
特别地它将给出所预告的配对 (\SGARefLink{sga7.sga7i.eix.n11-3-1}{11.3.1}). 考虑 $I$-模的正合列

\begin{equation}\tag{\SGARefTarget{sga7.sga7i.eix.n11-3-6}{11.3.6}}
0\longrightarrow U\xrightarrow{\ i\ }
 U\otimes_{\mathbb Z_\ell}\mathbb Q_\ell
 \xrightarrow{\ j\ }U\otimes D_\ell\longrightarrow0
\end{equation}
以及相应的上同调正合列, 这里 $I$ 的上同调用取值于所考虑拓扑群的\underline{连续上链}定义:

\begin{equation}\tag{\SGARefTarget{sga7.sga7i.eix.n11-3-7}{11.3.7}}
\resizebox{\textwidth}{!}{%
\begin{tikzcd}[ampersand replacement=\&, column sep=small, row sep=large]
0 \arrow[r]
 \& H^o(I,U) \arrow[r, "i^o"] \arrow[d, no head, "\wr"]
 \& H^o(I,U\otimes\mathbb Q_\ell) \arrow[r, "j^o"]
      \arrow[d, no head, "\wr"]
 \& H^o(I,U\otimes D_\ell) \arrow[r, "\partial"]
      \arrow[d, no head, "\wr"]
 \& H^1(I,U) \arrow[r, "i^1"]
 \& H^1(I,U\otimes\mathbb Q_\ell) \arrow[d, no head, "\wr"] \\
{}
 \& U^I
 \& (U\otimes\mathbb Q_\ell)^I\simeq U^I\otimes_{\mathbb Z}\mathbb Q
 \& (U\otimes D_\ell)^I
 \& {}
 \& H^1(I,U)\otimes_{\mathbb Z}\mathbb Q .
\end{tikzcd}}
\end{equation}
由 (\SGARefLink{sga7.sga7i.eix.n11-3-4}{11.3.4}) 定义的群 $\Phi$ 显然恰为 $\operatorname{coker}j^o$, 因而典范同构于 $\ker i^1$. 于是得到典范正合列

\begin{equation}\tag{\SGARefTarget{sga7.sga7i.eix.n11-3-8}{11.3.8}}
0\longrightarrow\Phi\longrightarrow H^1(I,U)
 \longrightarrow H^1(I,U)\otimes_{\mathbb Z}\mathbb Q,
\end{equation}
这表明 $\Phi$ \underline{典范同构于 $H^1(I,U)$ 的挠子群}.

另一方面, 由正合列 (\SGARefLink{sga7.sga7i.eix.n11-3-3}{11.3.3}), 对每个 $n=\ell^r$ 得到同构

\begin{equation}\tag{\SGARefTarget{sga7.sga7i.eix.n11-3-9}{11.3.9}}
H^1(I,T_n)\simeq\operatorname{Hom}(I,T_n)
 \xleftarrow{\sim}\operatorname{Hom}(T,T_n)
 \xleftarrow{\sim}\mathbb Z/n\mathbb Z.
\end{equation}
(注意 $I$ 在 $T$ 上作用平凡, 从而在 $T_n=T/nT$ 上也如此.) 因而对每个有限 $I$-$\mathbb Z_\ell$-模 $M$, 若以

\begin{equation}\tag{\SGARefTarget{sga7.sga7i.eix.n11-3-10}{11.3.10}}
M'=\operatorname{Hom}(M,T\otimes D_\ell)
\end{equation}
表示它的「Cartier 对偶」, 则存在典范配对

\begin{equation}\tag{\SGARefTarget{sga7.sga7i.eix.n11-3-11}{11.3.11}}
H^i(I,M)\otimes H^{1-i}(I,M')
 \longrightarrow H^1(I,T\otimes D_\ell)\simeq D_\ell,
\end{equation}
最后一个同构由同构 (\SGARefLink{sga7.sga7i.eix.n11-3-9}{11.3.9}) 对 $\ell$ 的各幂 $n$ 取归纳极限得到. 众所周知, 这些配对是\underline{完美对偶} (例如参见 SGA 5 I 5). 在这些配对中取 $M$ 形如 $U_n$, 从而 $M'=U'_n$, 再取极限, 得到完美配对

\begin{equation}\tag{\SGARefTarget{sga7.sga7i.eix.n11-3-12}{11.3.12}}
H^o(I,U\otimes D_\ell)\times H^1(I,U')\longrightarrow D_\ell,
\end{equation}
其中左边的第二个因子为有限生成 $\mathbb Z_\ell$-模. 再对配对 (\SGARefLink{sga7.sga7i.eix.n11-3-12}{11.3.12}), 沿形如 $\ell^r\operatorname{id}$ 的过渡态射取极限, 得到第二个完美配对

\begin{equation}\tag{\SGARefTarget{sga7.sga7i.eix.n11-3-13}{11.3.13}}
H^o(I,U\otimes\mathbb Q_\ell)
 \times H^1(I,U'\otimes\mathbb Q_\ell)
 \longrightarrow D_\ell
\end{equation}
它是 $\mathbb Q_\ell$ 上有限维向量空间的配对. 令 $i'$ 与 $j'$ 为如 (\SGARefLink{sga7.sga7i.eix.n11-3-6}{11.3.6}) 中定义的态射, 其中以 $U'$ 代替 $U$. 随即容易看出, 下列同态

\begin{center}
\begin{tikzcd}[column sep=6em, row sep=5em]
H^o(I,U\otimes\mathbb Q_\ell) \arrow[r, "j^o"]
  & H^o(I,U\otimes D_\ell) \\
H^1(I,U'\otimes\mathbb Q_\ell)
  & H^1(I,U') \arrow[l, "{i'}^1"']
\end{tikzcd}
\end{center}
关于配对 (\SGARefLink{sga7.sga7i.eix.n11-3-12}{11.3.12}) 与 (\SGARefLink{sga7.sga7i.eix.n11-3-13}{11.3.13}) 互为\underline{转置}. 立即推出, $\operatorname{coker}j^o$ 与 $\ker {i'}^1$ 通过一个取值于 $D_\ell$ 的典范配对互为对偶; 这就是所预告的配对 (\SGARefLink{sga7.sga7i.eix.n11-3-5}{11.3.5}).

% END INCLUDED SGA7I_English_current_buildable_source_20260823/components/167_expose_IX_section_11_3.tex

% BEGIN INCLUDED SGA7I_English_current_buildable_source_20260823/components/168_expose_IX_sections_11_4_and_11_5.tex
% Expose IX, sections 11.4 and 11.5.
% French transcription lines 3148--3191; authority scan physical pp. 466--468
% (scan indices 465--467; source folios 454--456). Physical p. 468 was
% checked through the exact 11.6 boundary. All pages were compared in four
% overlapping 1100-dpi bands. The text, formulas, primes, checks, underlines,
% and displayed arrows are clear at that resolution; no escalation was needed.

\medskip
\noindent
\SGARefTarget{sga7.sga7i.eix.n11-4}{11.4}. 上述利用连通分支群的 $\ell$-准素分量的表达式 (\SGARefLink{sga7.sga7i.eix.n11-2-1}{11.2.1}) 对配对 (\SGARefLink{sga7.sga7i.eix.n11-3-1}{11.3.1}) 的描述, 仅在 $\ell\ne p$ 时有效. 现在假设 $A_K$ 在 $S$ 上有\underline{半稳定约化}. 将再次定义为完美对偶的配对

\begin{equation}\tag{\SGARefTarget{sga7.sga7i.eix.n11-4-1}{11.4.1}}
\Phi_o(\ell)\otimes\Phi'_o(\ell)\longrightarrow\mathbb Q/\mathbb Z
\end{equation}
而不对 $\ell$ 作任何限制. 如同 \SGARefLink{sga7.sga7i.eix.n11-3}{11.3},\footnote{印本在 \SGARefLink{sga7.sga7i.eix.n11-4}{11.4} 本节内此处写作「如同 \SGARefLink{sga7.sga7i.eix.n11-4}{11.4}」. 此引用改为紧接在前的 \SGARefLink{sga7.sga7i.eix.n11-3}{11.3} 中的构造.} 应当证明这恢复了 (\SGARefLink{sga7.sga7i.eix.n1-2-1}{1.2.1}) 中定义的配对 (可能差一个应当确定的符号). 这样便能在半稳定约化情形建立猜想 \SGARefLink{sga7.sga7i.eix.n1-3}{1.3}, 如同限制到 $\ell\ne p$ 时的 $\ell$-准素分量一样.

为定义 (\SGARefLink{sga7.sga7i.eix.n11-4-1}{11.4.1}), 实质上将给出 $\Phi_o(\ell)$ 以 $K$ 上的 Barsotti--Tate 群 $T_\ell(A_K)$ 表示的表达式; 在当前情形, 它的作用与 \SGARefLink{sga7.sga7i.eix.n11-2}{11.2} 在上面处理的情形中的作用相同.

\medskip
\noindent
\underline{定理} \SGARefTarget{sga7.sga7i.eix.n11-5}{11.5}. \underline{假设 $S$ 为 Hensel 性状, 且 $A_K$ 在 $S$ 上有半稳定约化. 设 $\ell$ 为素数, 考虑单值配对} (\SGARefLink{sga7.sga7i.eix.n9-1-2}{9.1.2}), \underline{从而考虑典范同态}

\begin{equation*}
\widetilde u_\ell:\underline M_\ell
 \longrightarrow\underline{\check M}'_\ell.
\end{equation*}
\underline{置 $\Phi(\ell)=\operatorname{coker}\widetilde u_\ell$. 则存在典范同构}

\begin{equation}\tag{\SGARefTarget{sga7.sga7i.eix.n11-5-1}{11.5.1}}
\Phi_o(\ell)\xrightarrow{\ \sim\ }\Phi(\ell)\times_S s.
\end{equation}

此陈述立即给出完美配对 (\SGARefLink{sga7.sga7i.eix.n11-4-1}{11.4.1}) 的定义. 事实上, 将 \SGARefLink{sga7.sga7i.eix.n11-5}{11.5} 用于 $A'_K$ 以代替 $A_K$, 并使用 (\SGARefLink{sga7.sga7i.eix.n10-2-2}{10.2.2}), 还得到典范同构

\begin{equation*}
\Phi'_o(\ell)\xrightarrow{\ \sim\ }\Phi'(\ell)\times_S s,
\qquad
\Phi'(\ell)=\operatorname{coker}\bigl(
 {}^t\!\widetilde u_\ell:\underline M'_\ell
 \longrightarrow\underline{\check M}_\ell\bigr).
\end{equation*}

\medskip
\noindent
\underline{注} \SGARefTarget{sga7.sga7i.eix.n11-5-2}{11.5.2}. a) 定理 11.5 特别地蕴涵 $\widetilde u_\ell$ 为同源, 即 $u_\ell$ 是分离的, 包括 $\ell=p$ 的情形, 此时该断言不能由 \SGARefLink{sga7.sga7i.eix.n9-6}{9.6} 的定义立即得到. 注意, 此处没有使用 (深刻得多的) 定理 \SGARefLink{sga7.sga7i.eix.n10-4}{10.4}, 它将在下一节中证明.

\smallskip
\noindent
b) 利用 \SGARefLink{sga7.sga7i.eix.n10-4}{10.4}, 它给出单值配对

\begin{equation*}
u:\underline M\otimes\underline M'\longrightarrow\mathbb Z_S,
\end{equation*}
并以 $\Phi$ 表示对应同态的余核

\begin{equation*}
\widetilde u:\underline M\longrightarrow\underline{\check M}',
\end{equation*}
得到 \SGARefLink{sga7.sga7i.eix.n11-5}{11.5} 的一个与 $\ell$ 的选取无关的等价表述, 即典范同构

\begin{equation}\tag{\SGARefTarget{sga7.sga7i.eix.n11-5-3}{11.5.3}}
\Phi_o\xrightarrow{\ \sim\ }\Phi\times_S s.
\end{equation}

% END INCLUDED SGA7I_English_current_buildable_source_20260823/components/168_expose_IX_sections_11_4_and_11_5.tex

